\label{chap:p3-83-44-clasificacion_numero}


\label{chap:cobertura-taxonomia-numeros}
\index{nombre HMT!taxonomie}
\index{transducteur de Hensel!décalage complet}
\index{nombre rationnel!périodicité}

Le \cref{chap:numero-medicion} a défini un nombre HMT comme une section
compatible du système inverse de survie, et non comme une suite de
chiffres dépourvue de généalogie. Il restait cependant deux questions qu'il
convient de séparer avec précision. La première est une question de
\emph{représentation} : quelles valeurs réelles le transducteur
3-adique peut-il représenter avant d'imposer les filtres particuliers de clôture ? La seconde est
une question de \emph{sélection} : quelles sections de l'espace brut sont
choisies par APP, TRIT et TPK lorsque l'on exige simultanément leurs
compatibilités ? Cette section résout la première, formule la seconde comme
une restriction explicite du domaine et extrait de la carte de base mille une
taxonomie complète des valeurs visibles. Les outils de base sont
l'arithmétique \(p\)-adique, la dynamique symbolique et la caractérisation positionnelle
des rationnels \cite{Koblitz1984,LindMarcus1995,HardyWright2008}.

La distinction est structurelle. Un langage capable de représenter tout ruban ne
sélectionne pas de ce seul fait le ruban de \(\pi\), de \(e\), de \(\varphi\) ou
de \(\alpha\). À l'inverse, une loi de sélection ne peut pas être évaluée si
l'espace dans lequel elle agit n'a pas été construit auparavant. Représentation et sélection
sont deux théorèmes consécutifs, et non deux noms pour une même étape.

\section{Transducteur de Hensel non filtré et décalage complet sur les fibres ternaires}

Soit \(p\) un nombre premier, soit \(d\geq1\) et soit
\(A\in\operatorname{GL}_d(\mathbb Z)\). Pour un vecteur ligne
\(x_t\in\mathbb Z_p^d\), nous définissons
\begin{equation}
 y_t=x_tA,
 \qquad
 u_t=y_t\bmod p\in\mathbb F_p^d,
 \qquad
 x_{t+1}=\frac{y_t-u_t}{p}.
 \label{eq:transductor-hensel-general}
\end{equation}
Le bloc \(u_t\) est l'émission visible et \(x_{t+1}\) est la mémoire qui
subsiste après retrait de ce chiffre.

\begin{theorem}[Conjugaison de Hensel avec le décalage complet]
\label{thm:shift-hensel-completo}
Pour tout \(T\geq1\), l'application
\[
 \Psi_T:(\mathbb Z/p^T\mathbb Z)^d
 \longrightarrow(\mathbb F_p^d)^T,
 \qquad
 x_0\longmapsto(u_0,\ldots,u_{T-1}),
\]
est bijective. Son inverse est
\begin{equation}
 x_0\equiv
 \sum_{t=0}^{T-1}p^t u_tA^{-(t+1)}
 \pmod {p^T}.
 \label{eq:inversa-hensel-finita}
\end{equation}
En passant à la limite inverse, on obtient un homéomorphisme
\[
 \Psi:\mathbb Z_p^d\xrightarrow{\ \sim\ }
 (\mathbb F_p^d)^{\mathbb N}.
\]
Si \(S_A\) désigne la mise à jour \(x_t\mapsto x_{t+1}\) et
\(\sigma\) supprime le premier bloc d'un ruban, alors
\[
 \Psi\circ S_A=\sigma\circ\Psi.
\]
\end{theorem}

\begin{proof}
De \cref{eq:transductor-hensel-general}, on obtient
\[
 x_t=(u_t+px_{t+1})A^{-1}.
\]
L'itération de cette identité donne
\[
 x_0=
 \sum_{t=0}^{T-1}p^t u_tA^{-(t+1)}
 +p^Tx_TA^{-T}.
\]
Le dernier terme s'annule modulo \(p^T\), ce qui prouve qu'un mot de
\(T\) blocs détermine au plus une classe initiale. Réciproquement, étant donné
un mot quelconque, on fixe \(x_T=0\) et l'on reconstruit à rebours au moyen de
\(x_t=(u_t+px_{t+1})A^{-1}\). L'égalité
\(x_tA=u_t+px_{t+1}\) vérifie que la classe obtenue émet le mot
prescrit. Par conséquent, \(\Psi_T\) est bijective.

Les bijections sont compatibles avec la réduction modulo \(p^T\). Leurs limites
inverses produisent l'homéomorphisme \(\Psi\) ; de manière équivalente, la série de
\cref{eq:inversa-hensel-finita} converge dans \(\mathbb Z_p^d\). Appliquer
\(S_A\) retire exactement le bloc \(u_0\), ce qui correspond donc à
appliquer \(\sigma\) au ruban.
\end{proof}

Dans l'instance HMT, \(p=3\), \(d=6\) et la matrice publiée a pour
déterminant un. C'est pourquoi
\begin{equation}
 \mathbb Z_3^6\simeq(\mathbb F_3^6)^{\mathbb N},
 \qquad |\mathbb F_3^6|=3^6=729.
 \label{eq:hmt-shift-729}
\end{equation}
Chaque ruban de blocs de six trits possède une unique mémoire 3-adique initiale.
La matrice détermine la mise en coordonnées et la dynamique de lecture ;
l'universalité de l'alphabet démontre qu'elle n'exclut aucun ruban.

\begin{corollary}[Entropie et mesure de Haar]
La mise à jour de Hensel brute est conjuguée au décalage unilatéral
complet de
\(729\) symboles. Son entropie topologique est \(6\log3\). La mesure de Haar de
\(\mathbb Z_3^6\) est transportée vers la mesure produit uniforme sur les
blocs.
\end{corollary}

\begin{proof}
Il existe \(729^T=3^{6T}\) cylindres de longueur \(T\), et chaque classe modulo
\(3^T\) a une masse de Haar \(3^{-6T}\). La conjugaison du
\cref{thm:shift-hensel-completo} transporte les deux propriétés vers le
décalage.
\end{proof}

La conclusion probabiliste concerne l'espace brut muni de Haar. Elle
n'affirme pas que les sections publiées des constantes ont été échantillonnées
selon cette mesure ni qu'elles sont algorithmiquement aléatoires.

\section{Carte ternaire équilibrée des entiers}

La carte 3-adique brute et la carte archimédienne n'épuisent pas les types numériques
de la classification HMT. Pour l'opérateur entier, on utilise le mot vide ainsi que les
mots réduits
\[
 \mathscr W_{\mathrm{bal}}
 :=\{\varnothing\}\cup
 \{a_1\cdots a_m:
 a_j\in\{-1,0,+1\},\ a_1\ne0\}.
\]
Nous définissons récursivement
\[
 b(\varnothing)=0,
 \qquad
 b(wa)=3b(w)+a.
\]

\begin{theorem}[Représentation ternaire équilibrée unique]
\label{pf84:C03}
L'évaluation
\[
 b:\mathscr W_{\mathrm{bal}}\longrightarrow\ZZ
\]
est bijective. Les mots dont le premier trit est \(+1\) représentent les
entiers positifs et ceux qui commencent par \(-1\), les négatifs.
\end{theorem}

\begin{proof}
Pour \(n\ne0\), le dernier trit \(a\) est le représentant unique de
\(n\bmod3\) dans \(\{-1,0,+1\}\). Le parent
\[
 \frac{n-a}{3}
\]
a une valeur absolue strictement inférieure. L'itération se termine en \(0\), de
sorte que tout entier possède un développement fini. L'unicité du résidu
équilibré à chaque pas prouve l'unicité du mot. Le signe est
fixé par le premier trit non nul.
\end{proof}

\begin{remark}
Cette bijection est une carte arithmétique exacte. Elle n'identifie pas tous les
mots équilibrés à des sections survivantes TPK ni ne remplace la
généalogie, la retenue ou l'orientation d'une route enrichie.
\end{remark}

\section{Surjectivité de la carte archimédienne non filtrée sur \(\mathbb R\)}

Concaténons les coordonnées des blocs
\(u_t\in\mathbb F_3^6\) pour obtenir un ruban
\((a_n)_{n\geq1}\in\{0,1,2\}^{\mathbb N}\). Son évaluation ternaire est
\[
 \operatorname{ev}_3((a_n))=
 \sum_{n\geq1}\frac{a_n}{3^n}\in[0,1].
\]
Nous adoptons la convention canonique qui exclut la représentation avec une queue
ultimement constante de deux seulement lorsque le même point possède un autre développement
dans la carte. L'extrémité \(1\) conserve donc son unique représentation
fractionnaire \(0.222\ldots{}_3\).

\begin{theorem}[Couverture archimédienne brute]
\label{thm:cobertura-arquimediana-bruta}
L'application
\[
 P_{\rm raw}:\mathbb Z_3^6
 \xrightarrow{\Psi}(\mathbb F_3^6)^{\mathbb N}
 \xrightarrow{\operatorname{flat}}\{0,1,2\}^{\mathbb N}
 \xrightarrow{\operatorname{ev}_3}[0,1]
\]
est continue et surjective. Par conséquent,
\[
 \mathbb Z_3^6/\!\sim_{P_{\rm raw}}
 \ \cong\ [0,1]
\]
comme espace de valeurs, où \(x\sim_{P_{\rm raw}}y\) si leurs évaluations
coïncident.
\end{theorem}

\begin{proof}
Tout chiffre ternaire appartient à un unique bloc consécutif de six chiffres,
de sorte que \(\operatorname{flat}\) est une bijection d'espaces de rubans.
Toute \(r\in[0,1]\) possède un développement ternaire ; le
\cref{thm:shift-hensel-completo} fournit l'unique mémoire 3-adique qui
émet ses blocs. La continuité se vérifie par cylindres : fixer \(T\)
blocs fixe \(6T\) chiffres et enferme la valeur dans un intervalle de diamètre
\(3^{-6T}\). L'identification du quotient à l'image est la factorisation
canonique d'une application surjective par ses fibres.
\end{proof}

Une seule carte 3-adique est compacte et ne peut donc pas couvrir de manière
continue une droite non bornée. La globalisation correcte utilise une famille
dénombrable de cartes.

\begin{corollary}[Couverture archimédienne dénombrable de \(\mathbb R\)]
\label{cor:atlas-hmt-reales}
Soit
\[
 \widehat X_{\rm raw}=\mathbb Z\times\mathbb Z_3^6,
 \qquad
 \widehat P_{\rm raw}(m,x)=m+P_{\rm raw}(x).
\]
Alors \(\widehat P_{\rm raw}\) est surjective sur \(\mathbb R\).
En conséquence, tout réel admet au moins une généalogie de Hensel brute et
\[
 \widehat X_{\rm raw}/\!\sim_{\widehat P_{\rm raw}}
 \ \cong\ \mathbb R
\]
comme ensemble de valeurs.
\end{corollary}

\begin{proof}
Pour \(r\in\mathbb R\), prenons \(m=\lfloor r\rfloor\) et appliquons le
\cref{thm:cobertura-arquimediana-bruta} à \(r-m\in[0,1)\).
\end{proof}

Ce résultat précise l'affirmation selon laquelle \(\mathbb R\) est l'image
archimédienne quotient d'une
structure généalogique : la projection oublie la mémoire 3-adique et conserve
le point évalué. L'affirmation démontrée ici concerne la complétion
\emph{brute}. Si \(X_{\rm surv}\subseteq\mathbb Z_3^6\) désigne le sous-espace
qui satisfait en outre une famille concrète de clôtures APP--TRIT--TPK, alors
\[
 P_{\rm raw}(X_{\rm surv})=[0,1]
\]
est un théorème supplémentaire. Le transducteur garantit une capacité complète ; la
survie détermine quelle partie de cette capacité est réalisée par le système
restreint.

On ne déduit pas non plus ici que la somme et le produit natifs de TPK descendent à
tout le quotient global. On peut transporter l'arithmétique usuelle depuis
\(\mathbb R\), mais démontrer qu'elle coïncide avec des opérations construites avant
l'évaluation exige le théorème de descente correspondant.

Dans cette construction, les cylindres ne sont pas des infinitésimaux non nuls de
\(\mathbb R\). Chaque cylindre de profondeur finie a un diamètre positif,
et la suite des diamètres converge vers zéro. Deux sections qui coïncident jusqu'à
la profondeur \(T\) et diffèrent ensuite forment une \emph{perturbation de
queue} ; leur différence archimédienne est bornée par le diamètre du cylindre
commun. C'est la formulation exacte des ``cylindres qui tendent vers zéro''
dans la droite réelle ordinaire.

\section{Carte positionnelle en base \(10^3\) et classification des sections publiées}

La complétion cubique
\(1000=729+243+27+1\) rend naturel le regroupement des chiffres décimaux en blocs
\(b_n\in\{0,\ldots,999\}\). Pour un ruban de blocs, nous définissons
\[
 \operatorname{ev}_{1000}(b_1,b_2,\ldots)
 =\sum_{n\geq1}\frac{b_n}{1000^n}.
\]
Le développement canonique exclut une queue ultimement constante de blocs
\(999\). Cette convention résout la duplicité des extrémités de cylindre.

\begin{theorem}[Taxonomie visible en base mille]
\label{thm:taxonomia-base-mil}
Soit \(x\in[0,1)\) et soit \((b_n)\) son développement canonique en base mille.
Alors :
\begin{enumerate}
\item \(x\) possède un développement fini si et seulement si, écrit comme
fraction irréductible \(a/q\), on a \(q\mid1000^N\) pour un certain \(N\) ;
de manière équivalente, les seuls nombres premiers qui divisent \(q\) sont \(2\) et \(5\).
\item \(x\) est rationnel si et seulement si \((b_n)\) est ultimement périodique.
\item \(x\) est irrationnel si et seulement si \((b_n)\) n'est pas ultimement
périodique.
\end{enumerate}
\end{theorem}

\begin{proof}
Un développement qui se termine à la profondeur \(N\) a la forme
\(D/1000^N\) ; une fois réduit, son dénominateur ne contient que les nombres premiers \(2\) et
\(5\). Réciproquement, si \(q\mid1000^N\), alors \(a/q\) s'écrit avec
\(N\) blocs.

Si, après un préfixe de longueur \(r\), se répète une période de longueur
\(s\) dont l'entier en base mille est \(P\), la queue est une série géométrique :
\[
 \frac{P}{1000^r(1000^s-1)},
\]
et la valeur est rationnelle. À l'inverse, en développant \(a/q\), les restes de la
division successive ne peuvent prendre que \(q\) valeurs. Un reste nul produit un
développement fini ; la répétition d'un reste produit une période. La troisième
affirmation est la négation exacte de la seconde.
\end{proof}

Le mot \emph{fraction} désigne une présentation \(a/q\) d'un rationnel,
et non une espèce numérique supplémentaire. Il ne faut pas non plus confondre périodicité et
algébricité. Un réel est algébrique lorsqu'il annule un polynôme non nul de
\(\mathbb Q[X]\) ; il est transcendant lorsqu'il ne le fait pas. Cette seconde
classification est indépendante de la base utilisée pour écrire les chiffres.
Par exemple,
\[
 \varphi^2-\varphi-1=0
\]
fait de \(\varphi\) un irrationnel algébrique, tandis que la transcendance
de \(e\) et \(\pi\) découle de théorèmes classiques externes à HMT\@.

\section{Concaténation, réduction nonadique et retenue}

La compatibilité entre la carte cubique et la racine numérique est une identité
arithmétique, et non une ressemblance visuelle. Si \(u_j\in\{0,\ldots,999\}\) et
\[
 D_n=1000D_{n-1}+u_n,
 \qquad D_0=0,
\]
alors \(D_n\) est l'entier obtenu en concaténant les \(n\) premiers
blocs.

\begin{theorem}[Congruence nonadique de contrôle de la carte cubique]
\label{thm:checksum-nonadico-base-mil}
Pour tout \(n\geq1\),
\[
 D_n\equiv\sum_{j=1}^n u_j\pmod9.
\]
Par conséquent, concaténation en base mille, réduction modulo neuf et racine numérique
sont compatibles.
\end{theorem}

\begin{proof}
Comme \(1000\equiv1\pmod9\), la récurrence donne
\(D_n\equiv D_{n-1}+u_n\pmod9\). L'itération à partir de \(D_0=0\) produit la
formule.
\end{proof}

Dans la clôture dodécaphasique de \(\alpha\), l'équation de retenue par bloc
a la forme
\[
 P_m+E_m-\Phi_m-K_m+c_{m+1}=A_m+1000c_m.
\]
Réduite modulo neuf,
\[
 A_m\equiv
 P_m+E_m-\Phi_m-K_m+c_{m+1}-c_m
 \pmod9.
\]
Si les retenues extrêmes s'annulent, la somme télescopique élimine la mémoire
intermédiaire. Le résidu global de la sortie coïncide alors avec le résidu
de la combinaison des entrées. Cette loi explique pourquoi le registre nonadique
contrôle la concaténation décimale ; elle ne sélectionne pas à elle seule les blocs qui
doivent être concaténés.

\section{Coordonnée nonadique relevée et retour de phase}

La réduction modulo neuf ne doit pas être confondue avec l'identification des
entiers qui possèdent le même résidu. Soit
\[
 I_9=\{1,2,\ldots,9\}.
\]
Pour chaque entier positif \(n\), nous définissons sa \emph{phase nonadique} et son
\emph{indice de tour} par
\begin{equation}
 \rho_9(n)=1+((n-1)\bmod9),
 \qquad
 \kappa_9(n)=\frac{n-\rho_9(n)}9.
 \label{eq:coordenada-nonadica-levantada}
\end{equation}
La première coordonnée est la racine numérique écrite dans l'alphabet
\(I_9\), de sorte que la classe résiduelle zéro est représentée par la phase \(9\).
La seconde enregistre combien de tours complets ont précédé la phase
visible.

\begin{theorem}[Décomposition nonadique avec mémoire]
\label{thm:descomposicion-nonadica-memoria}
L'application
\[
 \Theta_9:\mathbb N_{>0}\longrightarrow\mathbb N_0\times I_9,
 \qquad
 n\longmapsto\bigl(\kappa_9(n),\rho_9(n)\bigr),
\]
est bijective, d’inverse
\[
 \Theta_9^{-1}(q,r)=9q+r.
\]
Dans ces coordonnées, le successeur arithmétique est conjugué à la dynamique
\begin{equation}
 \mathcal S_9(q,r)=
 \begin{cases}
  (q,r+1),&1\leq r<9,\\
  (q+1,1),&r=9.
 \end{cases}
 \label{eq:sucesor-nonadico-levantado}
\end{equation}
En particulier, le pas visible \(9\to1\) est un retour de phase accompagné
par l'incrément irréversible \(q\to q+1\) ; ce n'est pas une égalité entre les
entiers \(9\) et \(10\).
\end{theorem}

\begin{proof}
La division euclidienne de \(n-1\) par \(9\) fournit de manière unique
\(n-1=9q+s\), avec \(q\in\mathbb N_0\) et \(0\leq s<9\). En prenant
\(r=s+1\), on obtient \(n=9q+r\), avec \(r\in I_9\), ce qui démontre à la fois
l'existence, l'unicité et la formule inverse. Si \(r<9\), ajouter une unité
incrémente seulement \(r\). Si \(r=9\), l'égalité
\(9q+9+1=9(q+1)+1\) donne la seconde branche de
\cref{eq:sucesor-nonadico-levantado}.
\end{proof}

Ce relèvement sépare trois notions que le chiffre isolé mélange :
\begin{enumerate}
\item l'unité d'avancée, réalisée par le successeur ;
\item la phase visible \(r\), qui parcourt le cycle nonadique ;
\item la mémoire accumulée \(q\), qui distingue les retours d'une même phase.
\end{enumerate}
La projection \((q,r)\mapsto r\) oublie le nombre de tours et satisfait
\[
 \rho_9(n+9)=\rho_9(n),
 \qquad
 \kappa_9(n+9)=\kappa_9(n)+1.
\]
C'est le modèle arithmétique élémentaire du mécanisme qui, dans l'état TPK
enrichi, conserve en outre les quotients de Hensel, l'orientation, la signature et la donnée
de frontière. Un retour de mot visible peut ainsi coexister avec une
généalogie distincte.

Zéro n'appartient pas à \(I_9\) : on l'ajoute comme point de base séparé avant
d'appliquer \(\Theta_9\). Pour les entiers négatifs, on conserve le module
positif de la décomposition et l'on joint une orientation
\(\varepsilon\in\{-1,+1\}\). En conséquence, ni le signe ni zéro ne se
déduisent de la phase résiduelle. Ce typage évite de transformer une carte
cyclique en une identité arithmétique fausse.

\begin{theorem}[Entier orienté et signe généalogique]
\label{thm:entero-orientado-signo-genealogico}
L'application
\[
\Theta_9^\pm:
\mathbb Z\setminus\{0\}
\longrightarrow
\{-1,+1\}\times\mathbb N_0\times I_9,
\]
\[
n\longmapsto
\bigl(\operatorname{sgn}(n),
      \kappa_9(|n|),
      \rho_9(|n|)\bigr)
\]
est bijective, d’inverse
\[
(\varepsilon,q,r)\longmapsto\varepsilon(9q+r).
\]
La négation arithmétique agit par
\[
(\varepsilon,q,r)\longmapsto(-\varepsilon,q,r):
\]
inverse l'orientation et conserve la grandeur, la phase et la mémoire des tours.
\end{theorem}

\begin{proof}
Le \cref{thm:descomposicion-nonadica-memoria} fournit de manière unique
\(|n|=9q+r\), avec \(q\in\mathbb N_0\) et \(r\in I_9\). Le signe d'un entier
non nul est également unique. La formule inverse recompose \(n\), et remplacer
\(n\) par \(-n\) ne change que \(\varepsilon\).
\end{proof}

Ainsi, un négatif n'est ni une phase résiduelle spéciale ni une quantité
``impossible''. C'est la même grandeur nonadique transportée avec une orientation
généalogique opposée. Cette structure sera prolongée en mécanique dimensionnelle
par l'antisection particule--antiparticule.

Itérer \(\kappa_9\) produit une tour finie pour tout entier ordinaire ; dans
la complétion inverse, une suite compatible de quotients peut être
infinie. La différence entre les deux cas anticipe la taxonomie HMT : la valeur
visible répond au développement positionnel, tandis que le type généalogique
répond à l'évolution complète de ses mémoires.

\section{2 axes de classification : valeur visible et généalogie de transport}

La taxonomie précédente classe la valeur projetée. HMT ajoute un second
axe que l'écriture décimale ordinaire écarte. Pour une section
\(x\in\Omega_\infty^{\rm Arch}\), nous distinguons :
\[
 \begin{array}{rcl}
 \text{type visible}&=&
 \text{fini, périodique ou non périodique ; algébrique ou transcendant},\\
 \text{type généalogique}&=&
 \text{classe de retenue, orientation, fermeture et monodromie de }x.
 \end{array}
\]
Deux sections peuvent partager le même réel et différer selon le second axe. Un
réel rationnel possède un développement visible ultimement périodique, mais une
présentation enrichie redondante peut conserver une mémoire interne non
périodique que l'évaluation oublie. La périodicité de la \emph{valeur} se
rapporte toujours à son développement canonique ; la périodicité de la
\emph{généalogie} se rapporte à l'état enrichi de TPK.

Cette séparation fournit une définition concise :
\begin{equation}
 \boxed{
 \text{nombre HMT}
 =
 \text{section compatible avec mémoire};
 \qquad
 \text{valeur réelle}
 =
 \text{sa classe d'évaluation}.}
 \label{eq:numero-hmt-y-valor-real}
\end{equation}
Mesurer dans cette carte consiste à coupler la section au lecteur, à fixer un cylindre et à
conditionner les prolongements compatibles. Augmenter la précision signifie
rétrécir le cylindre, et non consommer la généalogie qui reste dans la fibre.

\section{Mémoire de profondeur arbitraire dans les systèmes d'émission finie}

La récurrence visible de neuf phases ne peut coexister avec une sortie
irrationnelle que si l'état complet conserve une information qui ne se réduit pas à
ces neuf étiquettes. Le résultat suivant rend cette nécessité exacte.

\begin{theorem}[Caractérisation par émetteurs finis]
\label{thm:emisor-finito-racional}
Soit \(Q\) un ensemble fini, \(F:Q\to Q\) une transition déterministe,
\(o:Q\to\{0,\ldots,999\}\) une sortie et \(q_0\in Q\). Le ruban
\(b_n=o(F^{n-1}(q_0))\) est ultimement périodique et son évaluation est rationnelle.
Réciproquement, tout ruban ultimement périodique admet un émetteur déterministe
fini.
\end{theorem}

\begin{proof}
Parmi \(|Q|+1\) états consécutifs, deux états égaux apparaissent. Dès la première
répétition, le déterminisme impose la répétition de tout le futur, de sorte que le
ruban est ultimement périodique. Le \cref{thm:taxonomia-base-mil} donne la
rationalité. Pour la réciproque, il suffit de construire une chaîne d'états pour le
préfixe et un cycle pour la période.
\end{proof}

Par conséquent, une loi qui engendre exactement \(e\), \(\pi\) ou tout autre
irrationnel ne peut pas consister uniquement en un automate déterministe à
états finis qui se répète sans mémoire supplémentaire. Elle peut utiliser une limite
inverse, une mémoire 3-adique, une profondeur croissante ou une entrée
apériodique ; tous ces mécanismes fournissent un état effectif non borné.
Dans HMT, le retour visible
\(g_{\rm vis}(n+9)=g_{\rm vis}(n)\) n'impose pas le retour de la retenue, de
l'orientation ni de la frontière. Cette différence est précisément l'endroit où
peut résider un développement non périodique.

\section{Le carré local de \texorpdfstring{\(e\)}{e} et la rupture de mémoire}

Pour un mot décimal \(a_1\cdots a_n\), nous appellerons \emph{carré local}
de longueur \(\ell\) à la position \(i\) une factorisation
\[
 a_i\cdots a_{i+2\ell-1}=vv,
 \qquad |v|=\ell.
\]
Cette notion relève de la combinatoire des mots. Elle n'affirme pas que la queue
infinie possède une période.

\begin{theorem}[Caractérisation finie du carré \(1828^2\)]
\label{thm:cuadrado-local-e}
Considérons les \(243\) mots du catalogue et le lecteur fixé
\(w_6\mapsto w_{30}\mapsto d_1d_2d_3d_4\), avec quatre triades en base mille.
Il existe un unique mot dont la lecture contient un carré de longueur quatre
commençant au deuxième chiffre décimal. C'est
\[
 w_e=201101,
 \qquad
 718281828459=7\,\underbrace{1828\,1828}_{(1828)^2}\,459.
\]
La répétition ne se prolonge pas en une troisième copie : le chiffre qu'exigerait cette
continuation serait \(1\), tandis que le chiffre effectif est \(4\). Dans tout le
catalogue, il n'existe qu'un autre carré de longueur quatre,
\[
 575157518472=\underbrace{5751\,5751}_{\text{carré local}}\,8472,
\]
qui commence au premier chiffre.
\end{theorem}

\begin{proof}
L'énumération exhaustive évalue les \(243\) mots par arithmétique
entière. Le profil des carrés pour les longueurs \(1,\ldots,6\) contient
\[
 240,\ 21,\ 3,\ 2,\ 0,\ 0
\]
occurrences, réparties entre
\[
 153,\ 19,\ 3,\ 2,\ 0,\ 0
\]
mots. Les deux occurrences de longueur quatre sont exactement celles de
l'énoncé. Deux énumérations entières indépendantes reproduisent le même recensement
et leurs sorties complètes coïncident terme à terme.
\end{proof}

Le mot \(w_e\) possède en outre la chaîne de relèvement
\[
201101\mid121221\mid102011\mid012222\mid102011.
\]
Les troisième et cinquième blocs visibles coïncident. En revanche,
les états qui les transportent ne coïncident pas. Avant les deux émissions, leurs réductions
modulo \(27\) sont
\[
(25,11,7,17,8,16),
\qquad
(7,8,4,23,26,7),
\]
et les quotients de Hensel ultérieurs sont
\[
(6,13,9,2,13,1),
\qquad
(15,0,3,19,23,25).
\]
Par conséquent, le retour \(102011\) est un retour de la projection visible, et non
de l'état enrichi. C'est la formulation exacte de la répétition locale
qui se rompt : le quotient conserve une généalogie que le bloc résiduel
n'enregistre pas.

En conséquence, le début
\[
e=2.718281828459\ldots
\]
ne présente pas de ``semi-périodicité'', mais un carré local de mot avec
retour visible et rupture de mémoire. Le résultat est exhaustif dans le cadre du
lecteur fixé ; son contenu est combinatoire et dynamique, et non une inférence
d'irrationalité à partir de douze chiffres.

\section[2 mesures finies de l'émetteur 468 à 243]{2 mesures finies de l'émetteur \texorpdfstring{\(468\to243\)}{468→243} : fréquence uniforme et priorité structurelle}

Le catalogue public distingue \(468\) émissions et \(104\,976\) germes.
Soit
\[
 q:E_{468}\longrightarrow W_{243}
\]
l'application qui oublie la signature d'émission et conserve le mot de six
trits. Il existe deux mesures naturelles, parce qu'il existe deux univers finis de dénombrement :
\[
 \mu_{468}(w)=\frac{|q^{-1}(w)|}{468},
 \qquad
 \mu_{\rm seed}(w)=
 \frac{\sum_{e\in q^{-1}(w)}m(e)}{104\,976},
\]
où \(m(e)\) est la multiplicité des germes de l'émission \(e\).

\begin{theorem}[Mesures canoniques relatives à l'univers dénombré]
\label{thm:medidas-emisor-468}
La mesure uniforme sur \(E_{468}\) est l'unique probabilité invariante sous
toutes les permutations de ses émissions ; sa mesure image est \(\mu_{468}\).
La mesure uniforme sur les \(104\,976\) germes induit
\(\mu_{\rm seed}\). Leurs histogrammes exacts sont
\[
\begin{array}{c|rrrrrr}
|q^{-1}(w)|&1&2&3&4&5&6\\ \hline
\#\{w\}&132&66&18&3&6&18
\end{array}
\]
et
\[
\begin{array}{c|rrrrrrrrr}
\text{poids des germes}
&144&288&432&576&864&1008&1440&1944&2088\\ \hline
\#\{w\}
&120&54&18&9&15&6&3&12&6.
\end{array}
\]
Pour les trois mots distingués par le lecteur, on obtient
\[
\begin{array}{c|cc}
&\mu_{468}&\mu_{\rm seed}\\ \hline
\pi&5/468&7/729\\
e&1/468&1/729\\
\varphi&1/468&1/729.
\end{array}
\]
\end{theorem}

\begin{proof}
L'invariance sous le groupe symétrique rend égales les masses de tous les
points de \(E_{468}\) ; la normalisation les fixe à \(1/468\). La formule de
la mesure image résulte de la sommation sur chaque fibre. Le même argument sur
l'ensemble des germes produit la seconde mesure. Les histogrammes s'obtiennent
par énumération exhaustive du catalogue fini de \(468\) émissions et de \(243\)
mots ; leurs sommes vérifient
\[
132+66+18+3+6+18=243,
\]
\[
132+2\cdot66+3\cdot18+4\cdot3+5\cdot6+6\cdot18=468
\]
et la somme totale des poids est \(104\,976\).
\end{proof}

La non-uniformité de la mesure image est un théorème fini : des mots distincts
ont des nombres distincts de réalisations. En revanche, il n'existe pas de
probabilité uniforme normalisable sur tout \(\mathbb R\), et ces mesures
n'ordonnent pas intrinsèquement tous les réels. Le mot de \(\pi\) appartient
à l'une des six fibres de taille cinq et de poids \(1008\) ; il n'est le maximum
unique d'aucun des deux histogrammes. Sa singularité démontrée procède
de la structure composée de ses cinq régions, de sa partition orientée
\(3+2\) et de ses compatibilités ultérieures, et non de la déclaration a priori d'une
probabilité absolue sur le continu.

\section{Sélecteur orbital fini antérieur au lecteur archimédien}
\label{sec:selector-orbital-constantes}

Les mesures précédentes révèlent des multiplicités, mais un mot isolé de
six trits dépend encore d'une orientation. Le catalogue possède une symétrie
interne qui permet de sélectionner d'abord des classes non orientées, puis de fixer
un représentant. Écrivons un mot comme \(abc\mid def\) et définissons
\[
 r(abc\mid def)=cab\mid efd,
 \qquad
 s(abc\mid def)=def\mid abc.
\]
Ces permutations agissent également sur les six coordonnées de la donnée
\(U_6\). Dans le catalogue complet, elles satisfont
\[
 r^3=1,
 \qquad s^2=1,
 \qquad srs=r^{-1},
\]
de sorte qu'elles engendrent une action que nous notons \(D_3^{\rm cat}\). L'action utilise
l'étiquetage TRIT fixe
\[
 \Psi(U_6)=(-U_6\bmod3)
\]
et la partition publiée des six coordonnées en deux triades.

Pour un mot \(w\), soient \(f(w)=|q^{-1}(w)|\) le nombre d'émissions,
\(m(w)\) sa masse de germes et
\[
 \nu(w)=(n_0(w),n_1(w),n_2(w))
\]
son recensement ordonné de trits. Soit en outre \(d(e)\) la classe diagonale additive d'une
émission et \(o(e)\in\{\mathrm{ES},\mathrm{NO}\}\) son orientation
publiée.

\begin{theorem}[Sélection orbitale interne des trois mots]
\label{thm:selector-orbital-constantes}
L'action \(D_3^{\rm cat}\) préserve l'ensemble des \(468\) émissions,
la projection \(q:E_{468}\to W_{243}\) et la multiplicité des germes. Son
action sur \(W_{243}\) possède exactement \(43\) orbites : \(38\) de taille
six et cinq de taille trois. Dans ce quotient, les unicités suivantes
sont vérifiées.

\begin{enumerate}
\item Il existe une unique orbite \(\mathcal O_\pi\) de taille six telle que, pour
tout \(w\in\mathcal O_\pi\),
\[
 f(w)=5,
 \qquad m(w)=1008,
 \qquad
 \{m(e):e\in q^{-1}(w)\}=\{432,144,144,144,144\}.
\]
C'est
\[
\mathcal O_\pi=
\{001112,010211,100121,112001,121100,211010\},
\]
et tous ses mots ont le recensement \(\nu=(2,3,1)\).

\item Considérons les orbites de taille six dont les mots ont une seule
émission, de multiplicité \(144\), et dont le support diagonal conjoint est
\(\{0,3,6\}\). Il y en a exactement six. Parmi elles, il existe une unique orbite
avec le même recensement \((2,3,1)\) que \(\mathcal O_\pi\) :
\[
\mathcal O_e=
\{011120,012110,101201,110012,120011,201101\}.
\]

\item Dans le même secteur, il existe une unique orbite équilibrée, c'est-à-dire avec
\(\nu=(2,2,2)\) :
\[
\mathcal O_\varphi=
\{002112,020211,112002,121200,200121,211020\}.
\]
\end{enumerate}

Enfin, la clause d'orientation
\begin{equation}
 \exists e\in q^{-1}(w):
 \qquad d(e)=6,
 \qquad o(e)=\mathrm{ES}
 \label{eq:gauge-orbital-constantes}
\end{equation}
sélectionne un unique représentant de chaque orbite et produit, respectivement,
\[
 \boxed{010211,\qquad201101,\qquad121200.}
\]
\end{theorem}

\begin{proof}
Les relations de groupe se vérifient directement comme des identités de
permutations de six coordonnées. L'énumération des \(468\) lignes
vérifie que \(r\) et \(s\) envoient chaque \(U_6\) sur une autre ligne du catalogue de
même multiplicité, et que la projection ternaire commute avec les deux
actions. Par conséquent, les orbites peuvent se calculer sur les \(243\) mots
sans perdre les fibres.

La partition exhaustive donne l'histogramme des tailles
\[
 38\cdot6+5\cdot3=243.
\]
En filtrant ces \(43\) orbites par le premier prédicat, il en reste exactement une ;
en filtrant par fibre à un point, multiplicité minimale et support diagonal
\(\{0,3,6\}\), il en reste six. Le recensement \((2,3,1)\) ne laisse qu'une orbite dans le
premier cas et le recensement équilibré \((2,2,2)\) n'en laisse qu'une dans le second.
Enfin, la recherche d'un relèvement avec \(d=6\) et orientation
\(\mathrm{ES}\) laisse un mot dans chacune des trois orbites. Toutes les
filtrations sont réalisées sur des descripteurs structurels du catalogue avant de consulter le
lecteur décimal.
\end{proof}

Les prédicats du théorème peuvent être réunis en un invariant de sélection.
Pour une orbite \(\mathcal O\), nous définissons son \emph{profil orbital}
\[
 \Sigma_{\rm cat}(\mathcal O)=
 \bigl(
 |\mathcal O|,\ f,\ m,\ \mathcal M_{\rm fib},\
 \mathcal D,\ \nu
 \bigr),
\]
où \(f\) est la taille de la fibre, \(m\) la masse des germes,
\(\mathcal M_{\rm fib}\) le multiensemble des multiplicités de ses
relèvements, \(\mathcal D\) le support diagonal lorsqu'il est employé et \(\nu\)
le recensement des trits. Les trois classes distinguées occupent des positions différentes
dans cet espace d'invariants :
\[
\begin{array}{c|c|c|c|c|c}
\mathcal O&|\mathcal O|&f&m&\mathcal M_{\rm fib}&(\mathcal D,\nu)\\ \hline
\mathcal O_\pi&6&5&1008&\{432,144,144,144,144\}&(-,(2,3,1))\\
\mathcal O_e&6&1&144&\{144\}&(\{0,3,6\},(2,3,1))\\
\mathcal O_\varphi&6&1&144&\{144\}&(\{0,3,6\},(2,2,2)).
\end{array}
\]
La singularité démontrée ne signifie pas que ces nombres reçoivent une
probabilité absolue sur \(\mathbb R\). Elle signifie qu'au sein de l'univers
fini produit par l'émetteur et avant l'évaluation des chiffres, leurs classes sont
séparables par des invariants internes et que la classe de \(\pi\) possède une fibre
régionale d'un type distinct. C'est la formulation mathématique d'une
priorité structurelle dans HMT\@.

Le contrôle négatif est informatif. Si l'on omet le support diagonal
\(\{0,3,6\}\), il existe quatre orbites à fibre unique équilibrées et
\(\mathcal O_\varphi\) cesse d'être unique. Avant de fixer la jauge, il y a six
orientations possibles dans chaque orbite. Après sa fixation, le mot de
\(\pi\) conserve trois relèvements positifs de \(U_6\), et non une seule région.
Ainsi, le théorème sélectionne le mot orienté mais respecte la multiplicité
géométrique des cinq régions.

La non-uniformité se renforce lors du passage aux orbites :
\[
\begin{array}{c|cc}
&\text{émissions}&\text{germes}\\ \hline
\mathcal O_\pi&5/78&14/243\\
\mathcal O_e&1/78&2/243\\
\mathcal O_\varphi&1/78&2/243.
\end{array}
\]
Ces quotients restent des masses du catalogue fini, et non des probabilités
attribuées aux nombres réels évalués.

L'avancée causale est précise. Les trois mots cessent d'être des entrées
nommées : ils sont retrouvés à partir de la symétrie et des invariants finis
d'APP--TPK, relativement à l'action \(D_3^{\rm cat}\), à l'étiquetage TRIT,
à la répartition \(3+3\) et à la jauge d'orientation de
\cref{eq:gauge-orbital-constantes}. La preuve n'utilise ni \(L_0\), ni \(L_1\),
ni chiffres décimaux, ni \(\alpha\), ni CODATA, ni l'état dodécaphasique. Restent des
affirmations ultérieures l'unicité absolue de cette action parmi tous les
automorphismes possibles, la déduction de la jauge au lieu de sa déclaration, la
canonicité du lecteur \(L_0/L_1\) indépendante de la valeur cible et l'émission coinductive des
développements infinis.

\section{Position de \texorpdfstring{\(\pi,e,\varphi\) et \(\alpha\)}{pi, e, phi et alpha} dans la classification}

Le \cref{thm:shift-hensel-completo} explique pourquoi tout préfixe
décimal peut être réémis à partir d'un état 3-adique : la capacité brute est
universelle. Par conséquent, la simple opération d'aller-retour sur un ruban ne sélectionne pas une
constante. Le contenu distinctif de la construction publiée se trouve
dans les objets antérieurs et postérieurs à cette réémission :
\begin{enumerate}
\item les fibres finies et les cinq régions distinctes de \(\pi\) ;
\item le lecteur archimédien typé qui associe mots, cylindres et blocs ;
\item la dodécaphase signée et la clôture réversible
\[
U_{12}^{\rm sgn}
\longleftrightarrow K
\longleftrightarrow(D_3K,D_4K,Q)
\longleftrightarrow\alpha_{\rm HMT};
\]
\item les critères de comparaison qui séparent une équivalence de cartes d'un
sélecteur antérieur à l'évaluation.
\end{enumerate}

Les douze triplets décimaux publiés de \(\pi,e,\varphi\) sont des coordonnées
finies exactes du lecteur enrichi. La chaîne imprimée de trente-six
chiffres est la projection dodécaphasique rationnelle
\(\alpha_{12}=\pi_{12}(\alpha_{\rm HMT})\) :
\[
\frac{7297352569283800997285105472380663}{10^{36}}.
\]
L'identité exacte dans cette fenêtre est
\[
\alpha_{12}
:=\pi_{12}(\alpha_{\rm HMT})
=\frac{7297352569283800997285105472380663}{10^{36}}.
\]
La récurrence compatible à profondeur arbitraire construit la section
infinie
\[
 \alpha_{\rm HMT}:=\varprojlim_N A^{(N)},
\]
et chaque fenêtre rationnelle est une projection exacte de ce même objet. Le
noyau \(E_{21/22}\) construit, dans une autre carte, la racine positive simple
finie \(\alpha_E\) ; le transport gradué du TPK prolonge ce noyau en
une famille compatible dont la racine limite est \(\alpha_B\). Si
\(\alpha_A:=\varprojlim_N\rho_N(\alpha_{12})\), les carrés de troncature
des deux familles commutent et déterminent à chaque profondeur le même
cylindre survivant. Le théorème exact des deux voies établit donc
\[
 \boxed{\alpha_A=\alpha_B=\alpha_{\rm HMT}}.
\]
L'égalité de \(\operatorname{Tr}_9(\alpha_E^{-1})\) et de
\(\operatorname{Tr}_9(\alpha_{12}^{-1})\) est le premier certificat fini de
cette identité ; elle ne la définit pas, ne limite pas sa portée et ne sélectionne pas les
prolongements. La projection typée conserve la distinction entre la section
coinductive, sa fenêtre rationnelle et la racine analytique sans séparer leur identité
à la limite.

La racine carrée de \(-1\) appartient à une autre carte. Dans le
\cref{chap:algebras-cuadraticas-orden-nueve}, on construit
\[
\mathbb F_9=\mathbb F_3[\iota]/(\iota^2+1),
\qquad \iota^2=-1,
\]
et l'on démontre que la matrice de Witt réalise cette structure quadratique. Il s'agit
d'une extension algébrique finie exacte, et non d'une catégorie visible de
développements réels.

\section{Carte de Hensel non filtrée, survie et frontière globale}

Les contrôles suivants sont des falsificateurs des résultats finis
précédents :
\begin{enumerate}
\item que deux états distincts modulo \(3^T\) émettent le même mot de
\(T\) blocs, ou qu'il existe un mot sans préimage ;
\item que l'inverse explicite de
\cref{eq:inversa-hensel-finita} échoue ;
\item qu'il apparaisse un rationnel dont le développement canonique ne soit pas ultimement
périodique, ou un irrationnel dont le développement le soit ;
\item qu'un émetteur déterministe fini produise un ruban non ultimement
périodique ;
\item que les recensements du catalogue ne totalisent pas \(468\), \(243\) et \(104\,976\),
ou que les mesures des trois mots ne coïncident pas avec l'énumération.
\end{enumerate}

La carte de Hensel brute
\[
\mathbb Z_3^6\longrightarrow[0,1]
\]
est continue et surjective ; en ajoutant la partie entière, son atlas brut couvre
\(\mathbb R\). Ce résultat exact n'implique pas que toute fibre brute contienne
un représentant qui survive aux filtres supplémentaires. Les fenêtres
finies publiées et les lecteurs nommés de \(\pi,e,\varphi\) sont vérifiés
dans HMT, relativement à leurs germes, régions,
orientations et règles enrichies.

\ifdefined\HMTInternalTraceability
Pour un préfixe filtré, soit \(T_u\) l'arbre des prolongements compatibles.
Si \(T_u\) est à ramification finie et possède un nœud à toute profondeur, le
lemme de König produit une branche infinie ; si, en outre, chaque niveau stable est
unitaire, la branche est unique. Cet énoncé est conditionnel : la non-vacuité universelle de \(T_u\)
pour tout préfixe n'a pas été démontrée. Restent donc
comme \textsc{programme ouvert} la couverture de toutes les fibres par des
survivants filtrés et la descente bien définie de la somme et du produit
natifs au quotient archimédien.
\fi




\label{chap:taxonomia-monodromia}

La représentation positionnelle d'un nombre enregistre une suite de symboles,
mais ne conserve pas nécessairement l'état qui détermine sa continuation.
L'holographie modulaire triadique distingue les deux niveaux. La phase nonadique indique
la position d'une lecture dans le cycle de neuf pas ; l'état
enrichi conserve, en outre, l'orientation, le quotient euclidien, le transport
d'entraînement et la condition de frontière. Le retour de la phase n'implique donc pas
le retour de l'état.

Soit \(S\) un espace d'états enrichis, soit
\(T:S\to S\) une transition et soit
\[
g:S\longrightarrow\mathbb Z/9\mathbb Z
\]
la phase, avec \(g(Ts)=g(s)+1\). Pour une base \(b\geq2\), un lecteur
\(d:S\to\{0,\ldots,b-1\}\) produit
\[
x(s)=\sum_{n\geq0}\frac{d(T^ns)}{b^{n+1}}.
\]
Lorsqu'une valeur admet les deux développements positionnels usuels, on adopte la
représentation canonique qui ne se termine pas par une queue constante \(b-1\). Cette
convention sépare une ambiguïté d'écriture de la multiplicité proprement
HMT, qui procède de sections enrichies distinctes ayant une même
évaluation.
L'opérateur de retour nonadique est la restriction
\[
M=T^9:g^{-1}(r)\longrightarrow g^{-1}(r).
\]
La différence entre \(T^9s=s\) et \(g(T^9s)=g(s)\) formalise la différence
entre périodicité et monodromie.

La transition de l'état enrichi de TPK peut être donnée par une relation et non par une
fonction. Dans ce cas, la notation précédente s'applique au décalage sur
l'espace des histoires infinies ou à une réalisation déterministe dont l'état
a été étendu avec l'information nécessaire pour fixer la continuation ; on ne
présuppose pas qu'un bloc visible possède un successeur unique.

\begin{definition}[Nombre HMT et présentation arithmétique]
Un nombre HMT est uniquement une section compatible
\(\boldsymbol x\) du système inverse d'états survivants. Un
lecteur positionnel \(L\) dont le domaine contient \(\boldsymbol x\) étant choisi, le couple
\((\boldsymbol x,L)\) est une présentation ou une mesure du nombre. La valeur
archimédienne publiée par cette présentation est l'intersection ponctuelle des
cylindres associés lorsque leurs diamètres tendent vers zéro. Deux sections ayant la
même valeur ne sont pas identifiées avant l'application du quotient archimédien.
\end{definition}

\begin{definition}[Périodicités visible et enrichie]
La sortie est éventuellement périodique lorsqu'il existe \(N\geq0\) et \(p>0\)
tels que
\[
d(T^{n+p}s)=d(T^ns),\qquad n\geq N.
\]
L'état est éventuellement périodique lorsque
\[
T^{N+p}s=T^Ns.
\]
La seconde condition implique la première. La réciproque exige que le lecteur
sépare les queues de l'orbite considérée, c'est-à-dire que
\[
 \bigl(d(T^{n+k}s)\bigr)_{k\geq0}
 =\bigl(d(T^{m+k}s)\bigr)_{k\geq0}
 \quad\Longrightarrow\quad T^ns=T^ms.
\]
\end{definition}

\begin{theorem}[Classification par retour de phase]
\label{thm:clasificacion-retorno-fase}
Soit \(s\in S\).
\begin{enumerate}
  \item Si l'orbite de \(s\) atteint un état dont la sortie future est
  constamment nulle, alors \(x(s)\) possède un développement fini et est
  rationnel.
  \item Si l'état est éventuellement périodique, alors la sortie est
  éventuellement périodique et \(x(s)\in\mathbb Q\).
  \item Si le lecteur sépare les queues de l'orbite et que l'orbite de \(s\) n'est pas
  éventuellement périodique, alors la sortie n'est pas éventuellement périodique
  et \(x(s)\notin\mathbb Q\).
  \item La non-périodicité ne distingue pas l’irrationalité algébrique de la
transcendance. Cette distinction exige un certificat polynomial ou un
théorème de transcendance supplémentaire.
\end{enumerate}
\end{theorem}

\begin{proof}
Les deux premiers points sont la caractérisation classique des rationnels
par des développements positionnels finis ou ultimement périodiques. Dans
le troisième, une périodicité ultime de la sortie produirait deux queues
égales. L’hypothèse de séparation identifierait alors les états
correspondants, ce qui contredirait la non-périodicité de l’orbite. Le quatrième
point découle de l’existence d’irrationnels algébriques et transcendants
dont les développements ne sont pas ultimement périodiques.
\end{proof}

\begin{corollary}[Monodromie nonadique et apériodicité]
Supposons que la phase revienne tous les neuf pas, que l’orbite de retour
\[
 \{M^ks:k\geq0\},\qquad M=T^9,
\]
ne soit pas ultimement périodique et que le lecteur sépare les queues. Alors la
sortie n’est pas ultimement périodique et le nombre produit est irrationnel. La
périodicité de la phase coexiste dans ce cas avec une évolution apériodique de
l’état enrichi.
\end{corollary}

L’hypothèse ne peut être remplacée par une seule inégalité \(M(s)\ne s\), ni
par un nombre fini de retours distincts. Relativement à l’état initial,
à la transition et au lecteur déclarés, la fenêtre actuelle certifie neuf
retours de phase avec changement d’état ; elle démontre une monodromie finie, mais non
à elle seule l’apériodicité de toute l’orbite.

\ifdefined\HMTRepetirResultadosTaxonomicos
\section{Le carré local de \texorpdfstring{\(e\)}{e} et la rupture de mémoire}

Pour un mot décimal \(a_1\cdots a_n\), nous appellerons \emph{carré local}
de longueur \(\ell\) à la position \(i\) une factorisation
\[
 a_i\cdots a_{i+2\ell-1}=vv,
 \qquad |v|=\ell.
\]
Cette notion appartient à la combinatoire des mots et ne doit pas être confondue avec
une période de la queue infinie.

\begin{theorem}[Caractérisation finie du carré \(1828^2\)]
Considérons les \(243\) mots du catalogue N33 et le lecteur figé
\(w_6\mapsto w_{30}\mapsto d_1d_2d_3d_4\), avec quatre triades en base mille.
Il existe un unique mot dont la lecture contient un carré de longueur quatre
commençant au deuxième chiffre décimal. C’est
\[
 w_e=201101,
 \qquad
 718281828459=7\,\underbrace{1828\,1828}_{(1828)^2}\,459.
\]
La répétition ne se prolonge pas en une troisième copie : le chiffre suivant requis
serait \(1\), tandis que le chiffre effectif est \(4\). Dans tout le catalogue,
il n’existe qu’un autre carré de longueur quatre,
\[
 575157518472=\underbrace{5751\,5751}_{\text{carré de mot}}\,8472,
\]
mais il commence au premier chiffre.
\end{theorem}

\begin{proof}
L’énumération exhaustive évalue les \(243\) mots par arithmétique
entière exacte. Le profil des carrés pour les longueurs \(1,\ldots,6\) contient,
respectivement,
\[
240, 21, 3, 2, 0, 0
\]
occurrences, réparties entre
\[
153, 19, 3, 2, 0, 0
\]
mots. Les deux occurrences de longueur quatre sont exactement celles écrites
dans l’énoncé. Deux énumérations entières indépendantes reproduisent le même
recensement et leurs sorties complètes coïncident terme à terme.
\end{proof}

Le mot \(w_e\) possède en outre la chaîne de relèvement
\[
201101\mid121221\mid102011\mid012222\mid102011.
\]
Le troisième et le cinquième blocs visibles coïncident. En revanche,
les états qui les transportent ne coïncident pas. Avant les deux émissions, leurs réductions
modulo \(27\) sont
\[
(25,11,7,17,8,16),
\qquad
(7,8,4,23,26,7),
\]
et les quotients de Hensel ultérieurs sont
\[
(6,13,9,2,13,1),
\qquad
(15,0,3,19,23,25).
\]
Par conséquent, le retour \(102011\) est un retour de la projection visible, non
de l’état enrichi. C’est la formulation exacte de l’intuition d’une
répétition locale qui se rompt : le quotient mémorise une généalogie que le
bloc résiduel ne conserve pas.

Le résultat appartient au lecteur enrichi déjà fixé par le système
APP--TRIT--TPK\@. Le recensement fini caractérise exhaustivement l’apparition du
carré local ; le prolongement coinductif de la section appartient à la
connexion nonadique construite dans le
\cref{chap:numero-heisenberg-d108}. La propriété combinatoire du préfixe et
la génération du développement complet sont donc deux niveaux
compatibles d’une même généalogie.

Cette formulation fournit le cadre exact pour interpréter le début
décimal de \(e\),
\[
e=2.718281828459\ldots.
\]
L’expression correcte n’est plus « semi-périodicité », mais \emph{carré local
de mot avec retour visible et rupture de mémoire}.

\section{2 mesures finies de l’émetteur \(468\to243\) : fréquence et priorité structurelle}

Il n’existe pas de distribution uniforme sur tous les nombres réels. Il
existe en revanche une mesure canonique sur le catalogue fini lorsque l’on munit d’une loi uniforme
l’ensemble des \(104\,976\) germes TPK\@. Pour un mot visible \(w\), on
définit
\[
\mu_{\mathrm{TPK}}(w)
=\frac{\#\{\text{germes se projetant sur }w\}}{104\,976}.
\]
La mesure n’est pas uniforme. Les \(243\) mots possèdent des poids compris
entre \(144\) et \(2088\), avec neuf valeurs de multiplicité distinctes.
L’histogramme exact, où le poids est écrit comme base et le nombre de mots
comme exposant, est
\[
144^{120},\ 288^{54},\ 432^{18},\ 576^9,\ 864^{15},\
1008^6,\ 1440^3,\ 1944^{12},\ 2088^6.
\]

Pour les trois mots distingués, on obtient
\[
\mu_{\mathrm{TPK}}(010211)=\frac{1008}{104\,976}=\frac7{729},
\]
\[
\mu_{\mathrm{TPK}}(201101)
=\mu_{\mathrm{TPK}}(121200)
=\frac{144}{104\,976}=\frac1{729}.
\]
Le mot associé à \(\pi\) a cinq antécédents \(U_6\) et un poids sept
fois supérieur à celui des mots associés à \(e\) et \(\varphi\). Ce résultat
démontre une hiérarchie combinatoire dans le catalogue ; il n’identifie pas à lui
seul les mots aux constantes classiques et n’établit pas de probabilité
sur \(\mathbb R\).

Il ne définit pas non plus, à lui seul, une priorité totale : six mots ont un poids
\(1008\), tandis que \(e\) et \(\varphi\) appartiennent à la classe minimale de
poids \(144\), partagée par \(118\) autres mots. La distinction HMT doit
s’exprimer par un profil structurel, par exemple
\[
\mathcal P(w)=
\bigl(m_{\rm semilla}(w),m_{U_6}(w),r_{\rm supervivencia}(w),
s_{\rm frontera}(w),s_{\rm excepcional}(w)\bigr),
\]
où chaque coordonnée est définie avant l’application d’étiquettes externes. Un scalaire
de « priorité » exige en outre de fixer, également à l’avance, la manière d’ordonner ou de
pondérer ces coordonnées.
\fi

\section{Classification des constantes selon l’antécédent généalogique, l’opération génératrice et le codomaine}

La classification HMT sépare des propriétés que la représentation ordinaire
réunit sous une seule valeur :
\begin{center}
\begin{tabular}{p{0.25\textwidth}p{0.64\textwidth}}
\toprule
Classe & Condition structurelle \\
\midrule
Fini & La sortie atteint une queue nulle. \\
Rationnel périodique & La sortie est ultimement périodique. \\
Rationnel monodromique & La sortie est périodique, mais l’état enrichi
parcourt une fibre non triviale que le lecteur visible oublie. \\
Irracional coinductivo & La sortie n’est pas ultimement périodique et les
cylindres déterminent une unique valeur. \\
Recurrente local & Un mot fini contient des retours ou des carrés, sans que
cela impose une périodicité à la queue. \\
Algébrique & La valeur satisfait un polynôme non nul à coefficients entiers ;
la route doit être accompagnée de ce certificat. \\
Transcendant & La valeur ne satisfait aucun polynôme entier non nul ; la
classification exige un théorème de transcendance. \\
Distinguido HMT & La section satisfait un prédicat structurel figé de
fibre, de frontière, de survie ou de symétrie ; c’est une classe transversale aux
précédentes. \\
\bottomrule
\end{tabular}
\end{center}

Ainsi, HMT ne modifie pas les définitions arithmétiques correctes de la rationalité,
de l’algébricité et de la transcendance. Elle les affine en ajoutant la généalogie de l’état,
la monodromie de phase et la multiplicité des fibres que l’évaluation
archimédienne avait oubliées.

En particulier, \(\varphi\) est algébrique parce qu’il satisfait
\(X^2-X-1=0\), tandis que la transcendance de \(e\) et de \(\pi\) découle de
théorèmes classiques externes. Ces propriétés sont compatibles avec le fait que les trois
sections soient distinguées par HMT\@. La clôture dodécaphasique construit plus
loin la projection rationnelle exacte
\(\alpha_{12}=\pi_{12}(\alpha_{\mathrm{HMT}})\), et le prolongement compatible
construit la section coinductive \(\alpha_{\mathrm{HMT}}\), dans le
\cref{chap:alpha-rev7} ; aucune d’elles n’est utilisée dans cette taxonomie comme entrée
ni publiée ici comme une seconde résidence. Le noyau \(E_{21/22}\)
produit une racine analytique distincte, dont la concordance avec la voie entière est
démontrée dans le codomaine de \(\operatorname{Tr}_9\). Section, projection
rationnelle, racine analytique, lecteur et valeur conservent ainsi leurs types propres.

Enfin, une « direction infinitésimale » exige un typage propre. Pour un
lecteur fini \(r_m\), deux sections distinctes dans une même fibre de \(r_m\)
définissent une différence invisible à la profondeur \(m\), quoique active dans un
prolongement ultérieur. Elles ne constituent pas pour autant un réel infinitésimal non nul.
Si deux sections coïncident sous toutes les troncatures fidèles du système
inverse, elles sont la même section. Un infinitésimal arithmétique véritable exigerait de
construire une extension ordonnée ou valuée non archimédienne et de démontrer que les
opérations HMT y descendent.



\section[Nombre comme section compatible]{Nombre comme section
compatible : classification arithmétique, évaluation archimédienne et
généalogie}
\label{p3-83-c44-s1:sec:ontologia-numeros-rev11}
\index{nombre!valeur et généalogie}
\index{taxonomie numérique}

La classification numérique de HMT conserve les espèces arithmétiques
ordinaires et ajoute une coordonnée que l’évaluation positionnelle écarte : la
généalogie de la section qui publie la valeur. Le résultat est une
classification conjointe. Le premier axe répond à la question du nombre obtenu ; le
second enregistre par quelle suite compatible d’états, de retours,
de retenues et de clôtures il est obtenu. Cette distinction permet de formuler dans un même
langage les entiers, les développements rationnels, les irrationnels, l’unité
imaginaire, les régularisations et les constantes spectrales sans identifier
des objets qui ne partagent qu’un scalaire.

\subsection{L’évaluation archimédienne sépare la section de sa valeur}

Soit
\[
 X_{\infty}^{\rm surv}
 =\varprojlim_n X_n^{\rm surv}
\]
un système inverse d’états survivants avec ses applications de
restriction. Une section compatible est une famille
\(\boldsymbol x=(x_n)_{n\geq0}\) qui satisfait
\(\rho_{n+1,n}(x_{n+1})=x_n\) à toute profondeur. Un lecteur positionnel
\(L\) associe à chaque préfixe un cylindre \(I_n(\boldsymbol x;L)\). Lorsque les
cylindres sont compatibles, que leurs diamètres convergent vers zéro et que la
convention canonique des extrémités est fixée, on définit
\begin{equation}
 \operatorname{Val}_{L}(\boldsymbol x)
 :=\text{le point de }
 \bigcap_{n\geq0}\overline{I_n(\boldsymbol x;L)}.
 \label{p3-83-c44-s1:eq:valor-lector-seccion-rev11}
\end{equation}

\HMTClaim{HMT-R-NUMBER-SECTION-VALUE}
\begin{definition}[Nombre, présentation et valeur dans HMT]
\label{p3-83-c44-s1:def:numero-presentacion-valor-rev11}
Un \emph{nombre HMT} est une section compatible
\(\boldsymbol x\in X_{\infty}^{\rm surv}\). La paire
\((\boldsymbol x,L)\) est une \emph{présentation} du nombre et
\(\operatorname{Val}_{L}(\boldsymbol x)\) est la valeur publiée par cette
présentation. La fibre
\[
 \operatorname{Val}_{L}^{-1}(r)
\]
conserve les généalogies que le quotient archimédien représente par
la même valeur \(r\).
\end{definition}

La capacité de la carte et la sélection d’une section sont des résultats
typés distincts. Le \cref{thm:cobertura-arquimediana-bruta} construit une
application continue et surjective
\[
 P_{\rm raw}:\mathbb Z_3^6\longrightarrow[0,1],
\]
et le \cref{cor:atlas-hmt-reales} la globalise en
\[
 \widehat P_{\rm raw}:\mathbb Z\times\mathbb Z_3^6
 \longrightarrow\mathbb R.
\]
Ces applications prouvent la capacité représentationnelle complète de la carte brute.
Le prédicat de survie et les lecteurs qui distinguent
\(\pi,e,\varphi\) sélectionnent, au sein de cette capacité, les généalogies
structurelles correspondantes.

\subsection{Terminaison et périodicité déterminent la rationalité}

Pour le développement canonique en base mille
\[
 x=\sum_{n\geq1}\frac{b_n}{1000^n},
 \qquad b_n\in\{0,\ldots,999\},
\]
on adopte la convention qui élimine la duplication produite par une queue
à partir d’un certain rang constante de blocs \(999\). Le
\cref{thm:taxonomia-base-mil} fournit alors la classification exacte
suivante.

\HMTClaim{HMT-R-NUMBER-DECIMAL-CLASSIFICATION}
\begin{theorem}[Classification visible dans la carte décimale cubique]
\label{p3-83-c44-s1:thm:clasificacion-visible-rev11}
Soit \(x=a/q\in[0,1)\) une fraction irréductible lorsque \(x\) est rationnel.
Alors :
\begin{enumerate}
\item le développement se termine si et seulement si \(q\mid1000^N\) pour un certain
      \(N\), équivalence qui restreint les facteurs premiers de \(q\) à
      \(2\) et \(5\) ;
\item le développement est périodique à partir d’un certain rang si et seulement si
      \(x\in\mathbb Q\) ;
\item le développement n’est pas périodique à partir d’un certain rang si et seulement si
      \(x\in\mathbb R\setminus\mathbb Q\).
\end{enumerate}
L’algébricité et la transcendance constituent un second axe : un réel est
algébrique s’il annule un polynôme non nul de \(\mathbb Q[X]\), et il est
transcendant s’il n’en annule aucun.
\end{theorem}

\begin{proof}
Un développement de longueur \(N\) a pour dénominateur un diviseur de \(1000^N\), et
la réciproque s’obtient en exprimant \(a/q\) avec ce dénominateur. Pour un
rationnel, la division successive possède un ensemble fini de restes : un reste
nul produit la terminaison et la répétition d’un reste produit la périodicité.
Une queue périodique est, réciproquement, une série géométrique rationnelle. La
classification polynomiale est indépendante de la base positionnelle.
\end{proof}

La dynamique enrichie distingue en outre la périodicité de sortie et la
périodicité d’état. Si \(T\) transporte l’état et \(d\) publie un
bloc, les conditions
\[
 d(T^{n+p}s)=d(T^n s),
 \qquad
 T^{n+p}s=T^n s
\]
définissent, respectivement, le retour visible et le retour de l'état. Le second
implique le premier. L'implication inverse exige que le lecteur sépare les
queues de l'orbite, conformément au
\cref{thm:clasificacion-retorno-fase}. C'est pourquoi une phase nonadique périodique
peut transporter une mémoire non périodique.

\begin{longtable}{@{}>{\raggedright\arraybackslash}p{.20\textwidth}
                    >{\raggedright\arraybackslash}p{.30\textwidth}
                    >{\raggedright\arraybackslash}p{.42\textwidth}@{}}
\caption{Taxonomie conjointe de la valeur publiée et de l'état enrichi.}
\label{p3-83-c44-s1:tab:taxonomia-conjunta-numeros-rev11}\\
\toprule
\textbf{Classe} & \textbf{Condition de la valeur} &
\textbf{Condition généalogique ou démonstration requise}\\
\midrule
\endfirsthead
\toprule
\textbf{Classe} & \textbf{Condition de la valeur} &
\textbf{Condition généalogique ou démonstration requise}\\
\midrule
\endhead
Entier & Élément de \(\mathbb Z\). & Mot ternaire équilibré fini et
unique ; le signe et la grandeur restent typés.\\
Decimal finito & Rationnel dont le dénominateur réduit divise une puissance
de \(1000\). & La sortie atteint une queue nulle.\\
Rationnel périodique & Développement ultimement périodique. & Un état
éventuellement périodique suffit ; une sortie périodique peut conserver une
fibre interne supplémentaire.\\
Rationnel monodromique & Même valeur rationnelle que la ligne précédente. & La
sortie revient tandis que l'état enrichi parcourt une fibre que le lecteur
visible oublie.\\
Irracional coinductivo & Développement non ultimement périodique. & Une famille
infinie de cylindres compatibles fixe une valeur unique ; l'apériodicité
globale exige une preuve sur toute l'orbite, et non une fenêtre finie.\\
Algébrique & Racine d'un polynôme non nul de \(\mathbb Q[X]\). & On joint
le polynôme et l'on vérifie son annulation.\\
Transcendant & Réel extérieur à la clôture algébrique de \(\mathbb Q\). & On
joint un théorème de transcendance ; la non-périodicité décimale à elle seule
ne démontre que l'irrationalité.\\
Distinguido HMT & L'une quelconque des classes arithmétiques précédentes. & La
section satisfait un prédicat figé de fibre, de frontière, de survie,
de symétrie ou de clôture.\\
\bottomrule
\end{longtable}

Ce tableau place sur des axes différents deux questions auxquelles il faut répondre
séparément. La classification ordinaire détermine les propriétés de la valeur ;
le lecteur HMT détermine la chaîne structurelle qui engendre la section. Par
exemple,
\[
 \varphi^2-\varphi-1=0
\]
classe \(\varphi\) comme irrationnel algébrique, tandis que la
transcendance de \(e\) et de \(\pi\) relève de théorèmes classiques externes.
Les trois peuvent, simultanément, être des sections distinguées par des prédicats
HMT différents. Leurs troncatures imprimées de mille ou dix mille chiffres sont
des rationnels finis qui démontrent les préfixes calculés d'une sortie prolongeable ; elles ne
remplacent pas l'objet infini qu'elles représentent.

La clôture dodécaphasique définit, quant à elle, la projection rationnelle finie
\begin{equation}
 \alpha_{12}=\pi_{12}(\alpha_{\rm HMT})
 =\frac{7297352569283800997285105472380663}{10^{36}}.
 \label{p3-83-c44-s1:eq:alpha-racional-taxonomia-rev11}
\end{equation}
Cette projection est exacte et rationnelle. L'objet
\(\alpha_{\rm HMT}=\varprojlim_NA^{(N)}\) est la section infinie construite
par la récurrence compatible ; ses troncatures rationnelles ne la remplacent pas.
La dérivation structurelle de la limite et la classification de chaque
projection finie sont des affirmations complémentaires.

\subsection{La carte ternaire équilibrée représente les entiers
bijectivement}

Soit
\[
 \mathscr W_{\rm bal}
 =\{\varnothing\}\cup
 \{a_1\cdots a_m:
 a_j\in\{-1,0,+1\},\ a_1\neq0\}
\]
et définissons
\[
 b(\varnothing)=0,
 \qquad
 b(wa)=3b(w)+a.
\]
Le \cref{pf84:C03} démontre que
\begin{equation}
 b:\mathscr W_{\rm bal}\xrightarrow{\ \sim\ }\mathbb Z
 \label{p3-83-c44-s1:eq:carta-entera-balanceada-rev11}
\end{equation}
est une bijection. Le dernier trit est le représentant unique de la
classe modulo trois et le quotient \((n-a)/3\) réduit strictement la valeur
absolue jusqu'à atteindre zéro. Le premier trit non nul fixe le signe. Cette carte
est une représentation arithmétique exacte ; l'appartenance d'un mot à une
section survivante incorpore en outre les conditions de route,
d'orientation et de frontière.

Pour \(n>0\), les coordonnées
\[
 \rho_9(n)=1+((n-1)\bmod9),
 \qquad
 \kappa_9(n)=\frac{n-\rho_9(n)}9
\]
réalisent la décomposition \(n=9\kappa_9(n)+\rho_9(n)\). La phase
\(\rho_9\in\{1,\ldots,9\}\) publie la classe résiduelle zéro comme \(9\), et
\(\kappa_9\) conserve le nombre de tours. Avec le signe
\(\varepsilon\in\{-1,+1\}\),
\[
 n\longleftrightarrow
 \bigl(\varepsilon,\kappa_9(|n|),\rho_9(|n|)\bigr)
\]
est bijectif sur \(\mathbb Z\setminus\{0\}\) ; zéro reste le point
de base. Le retour visible \(9\to1\) est ainsi une transition de phase avec
incrément de mémoire, et non une identification d'entiers consécutifs.

\subsection{Structure complexe finie induite par l'endomorphisme de Witt \texorpdfstring{\(J^2=-I_6\)}{J²=-I₆}}

L'unité imaginaire ordinaire est l'élément
\(i\in\mathbb C\) qui satisfait \(i^2=-1\) ; il est algébrique de degré deux sur
\(\mathbb R\) et sur \(\mathbb Q\). La réalisation HMT appartient à un
domaine fini différent. Si \(C_P\) est la matrice de conférence de Paley et
\(J=A_W\) sa réduction modulo trois dans
\(\operatorname{End}_{\mathbb F_3}(\mathbb F_3^6)\), alors
\begin{equation}
 \boxed{J^2=2I_6=-I_6.}
 \label{p3-83-c44-s1:eq:raiz-interna-menos-uno-ontologia-rev11}
\end{equation}
Comme \(X^2+1\) est irréductible dans \(\mathbb F_3[X]\), l'extension
\[
 \mathbb F_9=\mathbb F_3[\iota]/(\iota^2+1)
\]
sépare les espaces propres de \(J\) associés à \(\pm\iota\). Le résultat,
démontré dans \cref{sec:raiz-interna-menos-uno-rev6}, dote le module ternaire d'une
structure complexe finie et organise deux orientations conjuguées. Sa force
est \emph{exacte interne, démontrée avec la matrice de Paley--Witt déclarée} ; la
comparaison avec \(\mathbb C\) conserve la distinction entre les deux corps.

\subsection{\(6\) opérateurs réalisent le rationnel
\texorpdfstring{\( -1/12 \)}{-1/12} dans des domaines distincts}

La régularisation zêta attribue à la suite des entiers positifs la partie
finie obtenue par prolongement méromorphe :
\begin{equation}
 \sum_{n\geq1}^{[\zeta]}n
 :=\zeta(-1)=-\frac1{12}.
 \label{p3-83-c44-s1:eq:regularizacion-zeta-menos-doce-rev11}
\end{equation}
L'exposant \([\zeta]\) fait partie du type de l'expression. La somme
partielle ordinaire satisfait
\(\sum_{n=1}^{N}n=N(N+1)/2\) et croît sans borne ; c'est pourquoi
\eqref{p3-83-c44-s1:eq:regularizacion-zeta-menos-doce-rev11} est une valeur régularisée, et non
la limite de ces sommes partielles.

HMT contient en outre six réalisations algébriques exactes du même
rationnel. La coïncidence scalaire permet de les comparer ; chaque objet conserve son
domaine, son application et sa normalisation.

\begin{longtable}{@{}>{\raggedright\arraybackslash}p{.23\textwidth}
                    >{\raggedright\arraybackslash}p{.43\textwidth}
                    >{\raggedright\arraybackslash}p{.26\textwidth}@{}}
\caption{Réalisations typées de \(-1/12\).}
\label{p3-83-c44-s1:tab:realizaciones-menos-doce-rev11}\\
\toprule
\textbf{Domaine} & \textbf{Application et évaluation} &
\textbf{Conclusion et normalisation}\\
\midrule
\endfirsthead
\toprule
\textbf{Domaine} & \textbf{Application et évaluation} &
\textbf{Conclusion et normalisation}\\
\midrule
\endhead
Incidence marquée hexade--octade &
\(F_6\hookrightarrow F_8\), avec
\(\lvert F_6\rvert=90\), \(\lvert F_8\rvert=120\), et
\[
 (90-120)/(24\binom62)=-1/12.
\]
& Défaut combinatoire exact pour la normalisation déclarée.\\
Périodicités dodécaphasiques &
Un pas de \(30^\circ\) produit
\[
 \begin{aligned}
 \frac{30}{120}-\frac{30}{90}
 &=\frac14-\frac13,\\
 &=-\frac1{12}.
 \end{aligned}
\]
& Différence orientée exacte de périodes.\\
Opérateur de Gram &
Pour \(K_{60,81}=\operatorname{diag}(0,12)\), la pseudo-inverse sur le
rang non nul vaut \(1/12\) ; l'orientation négative publie \(-1/12\). &
Identité opératorielle exacte sur la composante non nulle.\\
Deuxième différence d'échelle &
Le calcul produit \(C(L)=8/81\) et l'extracteur déclaré
\(R(L)=-(27/32)C(L)=-1/12\). & Exact après fixation du facteur de
normalisation \(-27/32\) ; ce facteur est une primitive explicite.\\
Quotient modulaire êta &
\[
 \begin{aligned}
 \frac{\eta(\tau)}{\eta(3\tau)}
 &=q^{-1/12}\\[-1mm]
 &\quad{}\times
 \prod_{n\geq1}(1+q^n+q^{2n})^{-1}.
 \end{aligned}
\]
& Exposant modulaire exact, provenant de \(1/24-3/24\).\\
Projecteurs du cycle de douze phases &
Pour \(D_r=I-S^r\),
\[
 \begin{aligned}
 \tau(\Pi_{\ker D_3}-\Pi_{\ker D_4})
 &=\frac{3-4}{12},\\
 &=-\frac1{12}.
 \end{aligned}
\]
& Défaut transversal exact entre les pas trois et quatre.\\
\bottomrule
\end{longtable}

Les six réalisations sont démontrées dans
\cref{sec:menos-doce-realizaciones-rev6}. La régularisation zêta, le défaut
d'incidence, la pseudo-inverse, l'extracteur d'échelle, le quotient êta et la
trace des projecteurs sont des applications différentes. Une identification de leurs
domaines nécessiterait un entrelaceur supplémentaire. De même, le terme
\(+1/12\) de la formule de Glaisher--Kinkelin relève de la régularisation
de l'hyperfactorielle et conserve ce type propre.

\subsection{Décomposition multiplicative des nombres premiers comme modes irréductibles du TPK}
\index{nombre premier!irréductible}
\index{horloge première}

Un nombre premier ordinaire est un entier \(p>1\) dont les factorisations
\(p=ab\), avec \(a,b\in\mathbb N\), contiennent nécessairement une unité. Le
coproduit réduit
\[
 \Delta_{\times}^{\rm red}e_n
 =\sum_{\substack{ab=n\\a,b>1}}e_a\otimes e_b
\]
convertit cette définition en un sélecteur spectral. Selon le
\cref{pf84:C07}, l'opérateur positif
\[
 N_{\times}^{\rm red}
 =(\overline{\Delta_{\times}^{\rm red}})^*
   \overline{\Delta_{\times}^{\rm red}}
\]
satisfait
\(N_{\times}^{\rm red}e_n=d_{\rm red}(n)e_n\), et son noyau, après
retrait de \(e_1\), sélectionne exactement les nombres premiers. L'entrelaceur
multiplicatif transporte le même sélecteur au produit natif d'APP\@. Ce
sélecteur est la projection spectrale sur
\(\ker N_{\times}^{\rm red}\ominus\mathbb C e_1\), qui coïncide exactement
avec le sous-espace engendré par les nombres premiers et ne dépend pas de la reconnaissance de listes
décimales.

La factorisation unique fournit en outre la réalisation de Fock
\[
 U_F:\Gamma_s(\ell^2(\mathbb P))
 \xrightarrow{\ \sim\ }\ell^2(\mathbb N_{\geq1}),
 \qquad
 U_F\,d\Gamma(H_{\mathbb P})\,U_F^*=\log Q,
\]
où \(H_{\mathbb P}e_p=(\log p)e_p\). Chaque nombre premier est un mode
multiplicativement irréductible ; un entier composé est une occupation finie
de ces modes.

Pour \(n\) premier avec \(10\), son horloge décimale est l'ordre multiplicatif
\[
 \tau_n=\operatorname{ord}_n(10).
\]
Si \(p\) est premier et \(p\notin\{2,3,5\}\), le
\cref{thm:cierre-nonadico-reloj-primo-rev6} démontre
\begin{equation}
 \tau_p\mid p-1,
 \qquad
 M_p:=\frac{10^{\tau_p}-1}{p}\in\mathbb N,
 \qquad
 M_p\equiv0\pmod9.
 \label{p3-83-c44-s1:eq:reloj-primo-cierre-nonadico-rev11}
\end{equation}
Pour un nombre composé,
\[
 \operatorname{ord}_{n}(10)
 =\operatorname{lcm}_{p^a\parallel n}\operatorname{ord}_{p^a}(10).
\]
Par conséquent, l'horloge composée synchronise des horloges de puissances premières, tandis
que le nombre premier apporte une composante irréductible. La congruence nonadique
\eqref{p3-83-c44-s1:eq:reloj-primo-cierre-nonadico-rev11} décrit la clôture de phase ; le
sélecteur de primalité réside dans le coproduit réduit. La séparation de
ces deux fonctions évite d'attribuer une capacité sélective à la seule période décimale.

\subsection{Chaque constante est associée à son opérateur d'extraction}

Les cycles d'ordre \(3^m\), la limite de leurs traces de chaleur et la
transformée thêta--Mellin construisent la fonctionnelle
\begin{equation}
 \mathcal Z_3(s)
 =\frac{\displaystyle
 \int_0^\infty\Theta_+(x)x^{s/2-1}\,\mathrm dx}
 {\pi^{-s/2}\Gamma(s/2)}
 =\zeta(s),
 \qquad \Re s>1.
 \label{p3-83-c44-s1:eq:funcional-triadico-resumen-rev11}
\end{equation}
L'inversion thêta prolonge méromorphiquement la fonctionnelle. La partition
résiduelle
\[
 \zeta(s)=Z_0(s)+Z_1(s)+Z_2(s)
\]
sépare agrégat, échelle et orientation. Une seule source opératorielle admet
plusieurs applications d'extraction ---partie finie, coefficient de Laurent,
évaluation, dérivée régularisée ou caractère--- ; la coïncidence de source
n'identifie pas les constantes résultantes.

% REMATE_LOCAL_CONSTANTES_TABLA
\par

\begingroup
\small
\fontsize{10pt}{11.2pt}\selectfont
\setlength{\LTpre}{0.5\baselineskip}
\setlength{\LTpost}{0.5\baselineskip}
\begin{longtable}{@{}>{\raggedright\arraybackslash}p{.15\textwidth}
                    >{\raggedright\arraybackslash}p{.245\textwidth}
                    >{\raggedright\arraybackslash}p{.305\textwidth}
                    >{\raggedright\arraybackslash}p{.195\textwidth}@{}}
\caption{Définition, application d'extraction et portée mathématique des
constantes numériques complémentaires.}
\label{p3-83-c44-s1:tab:constantes-operadores-ontologia-rev11}\\
\toprule
\textbf{Objet} & \textbf{Définition standard} &
\textbf{Application et référence mathématique} & \textbf{Conclusion et condition}\\
\midrule
\endfirsthead
\toprule
\textbf{Objet} & \textbf{Définition standard} &
\textbf{Application et référence mathématique} & \textbf{Conclusion et condition}\\
\midrule
\endhead
Euler--Mascheroni \(\gamma\) &
Pour \(H_N=\sum_{n\leq N}n^{-1}\),
\(\gamma=\lim_{N\to\infty}(H_N-\log N)\), équivalente à la partie finie de
\(\zeta\) en \(s=1\). &
Si \(C_r=\operatorname*{FP}_{s=1}Z_r(s)\), alors
\(\gamma=C_0+C_1+C_2\), par
\eqref{eq:residual-three-covectors}. &
Identité exacte de la fonctionnelle résiduelle démontrée avec l'opérateur et la
normalisation déclarés.\\
Constantes de Stieltjes \(\gamma_n\) &
Si \(\zeta(1+z)-z^{-1}=\sum_{n\geq0}a_nz^n\), elles sont définies par
\(\gamma_n=(-1)^n n!a_n\). &
Les coefficients résiduels \(c_{r,n}\) satisfont
\(\gamma_n=(-1)^n n!\sum_r c_{r,n}\) et
\(L^{(n)}(1,\chi_{-3})/n!=c_{1,n}-c_{2,n}\) : c'est la décomposition
coefficient par coefficient de \eqref{eq:stieltjes} selon les trois classes
résiduelles de \eqref{eq:residual-three-covectors}. &
Décomposition exacte ; les extractions numériques vérifient les
coefficients, et non leur définition.\\
Ap\'ery \(\zeta(3)\) &
Valeur de la zêta de Riemann en \(s=3\). &
\(\mathcal Z_3(3)=\zeta(3)\), avec des préfixes bornés par des intervalles rationnels ;
sa résidence spectrale dans la tour graduée est fixée dans
\cref{eq:bendersky,eq:odd-zeta-bendersky}. &
Évaluation exacte de la fonctionnelle ; son irrationalité découle du théorème
classique d'Apéry \cite{Apery1979}.\\
Glaisher--Kinkelin \(A_{\rm GK}\) &
Constante de la croissance régularisée de l'hyperfactorielle. &
\(A_{\rm GK}=\exp(1/12-\mathcal Z_3'(-1))\), par \cref{eq:bendersky,eq:odd-zeta-bendersky}. &
Dérivée régularisée exacte ; \(1/12\) relève de cette normalisation
hyperfactorielle.\\
Euler--Kronecker \(\gamma_{F_{-3}}\) &
Partie finie logarithmique de la zêta de Dedekind du corps
\(F_{-3}=\mathbb Q(\sqrt{-3})\). &
Le caractère résiduel \(\chi_{-3}\) sélectionne le corps et produit
\(\gamma_{F_{-3}}=\gamma+L'(1,\chi_{-3})/L(1,\chi_{-3})\), par
\eqref{eq:euler-kronecker}. &
Identité exacte ; le choix du corps procède du caractère modulo trois.\\
Meissel--Mertens \(B_1\) &
Pour \(S_P(x)=\sum_{p\leq x}p^{-1}\),
\(B_1=\lim_{x\to\infty}(S_P(x)-\log\log x)\). &
C'est la partie finie de la trace première tronquée sur le secteur à une particule
du hamiltonien de Fock, par \eqref{eq:meissel-mertens} ; le sélecteur de
modes premiers est fixé avant d'effectuer la régularisation. &
Représentation opératorielle exacte une fois les modes premiers fixés ; le
sélecteur d'irréductibles est antérieur.\\
Feigenbaum \((\delta,\alpha_{\rm F})\) &
Constantes universelles d'échelle du point fixe de renormalisation pour la
cascade quadratique de doublement de période. &
La famille dodécaphasique déclarée produit au niveau huit
\(\delta^{\mathscr B}_8(0)=4.66921218915\ldots\) et
\(\alpha^{\mathscr B}_8(0)=2.50296847988\ldots\), par
\cref{sec:p3-final:cierre-feigenbaum}. &
Calcul fini exact jusqu'à \(n=8\) ; la convergence vers le point fixe universel
exige de démontrer la convergence de l'opérateur de renormalisation.
\(\alpha_{\rm F}\) et \(\alpha_{\rm HMT}\) sont
des objets distincts.\\
Catalan \(G\) &
\(G=L(2,\chi_{-4})=\sum_{n\geq0}(-1)^n/(2n+1)^2\). &
Le lecteur naturel est le caractère impair modulo quatre évalué en \(2\),
par \cref{hmtctx:catalan:sec:combinaciones-lectores-vecinos-rev6}. &
Identité classique exacte avec opérateur localisé ; ce n'est pas une sortie de la
fonctionnelle résiduelle modulo trois.\\
Khinchin \(K\) &
Moyenne géométrique presque sûre des chiffres de fraction continue pour le
système de Gauss. &
Le lecteur est la moyenne ergodique de \(\log a_1\) pour
\(T(x)=x^{-1}-\lfloor x^{-1}\rfloor\) et sa mesure invariante, par
\cref{hmtctx:khinchin:sec:combinaciones-lectores-vecinos-rev6}. &
Théorème ergodique classique incorporé ; il exige une dynamique différente de
la tour spectrale triadique.\\
\bottomrule
\end{longtable}
\endgroup

Le tableau fixe une règle de lecture : une constante est définie par son
objet, son opérateur et son application d'extraction, et non par l'apparition d'un
préfixe décimal dans un catalogue. Euler--Mascheroni et Stieltjes relèvent des
parties finies et des jets de Laurent ; Apéry, d'une évaluation ;
Glaisher--Kinkelin, d'une dérivée régularisée ; Euler--Kronecker, du caractère
du corps ; Meissel--Mertens, d'une trace première ; Feigenbaum, d'une famille
dynamique finie ; Catalan et Khinchin, d'opérateurs voisins distincts.

\subsection{Domaines, prémisses et conditions nécessaires de la classification généalogique du nombre}

Les définitions arithmétiques d'entier, de rationnel, d'irrationnel, d'algébrique et de
transcendant sont standard. La carte ternaire équilibrée, la séparation entre
section et valeur, le sélecteur natif d'irréductibles, la racine de Paley--Witt de
moins un, les six réalisations de \(-1/12\) et les lecteurs spectraux se
déduisent des applications citées dans chaque partie. Cette section les réunit sans
introduire de constante ni de sélecteur numérique supplémentaire.

La force de chaque affirmation est déterminée par son application :
\begin{enumerate}
\item les bijections, identités matricielles, coproduits, traces et
      évaluations indiqués sont exacts dans leurs domaines déclarés ;
\item la transcendance de \(e\) et \(\pi\), l'irrationalité de
      \(\zeta(3)\) et le théorème ergodique de Khinchin sont des théorèmes classiques
      incorporés comme contrôle ou classification ;
\item les approximants de Feigenbaum démontrent la branche et la profondeur
      calculées ; l'universalité exige l'opérateur de renormalisation et
      son théorème de convergence ;
\item l'égalité de deux scalaires provenant de domaines différents
      autorise une comparaison et exige un entrelaceur pour identifier les
      objets complets ;
\item une troncature décimale démontre le préfixe calculé, tandis que la
      classe arithmétique de la valeur infinie exige, selon le cas, une équation
      polynomiale ou un théorème d'irrationalité ou de transcendance.
\end{enumerate}

Avec ces règles, l'ontologie HMT du nombre acquiert une formulation
précise : la valeur archimédienne est la coordonnée publiée ; la section
compatible conserve la chaîne qui l'engendre ; l'opérateur détermine quelle
propriété est lue ; et ses hypothèses déterminent exactement quelle conclusion s'en
déduit.


\section[Clôture généalogique du concept de nombre]{Clôture généalogique du concept de nombre : compatibilité, mémoire et évaluation archimédienne}
\label{p3-83-c44-s2:chap:cierre-genealogico-numero}
\index{nombre!clôture généalogique}
\index{chiffre!coordonnée visible}
\index{valeur!évaluation archimédienne}
\index{conservation évolutive de l'unité}

La théorie des nombres développée dans HMT culmine dans une distinction de types
qui gouverne la transition ultérieure vers la mécanique dimensionnelle. Un
\emph{chiffre} est une coordonnée visible d'un état fini ; une \emph{valeur}
est l'image archimédienne publiée par un lecteur ; un \emph{nombre HMT} est
la généalogie compatible complète qui conserve, avec ses préfixes,
orientation, quotient, retenue, frontière et mémoire. La fibre d'une valeur
réunit les généalogies distinctes qui publient ce même scalaire.

Cette construction élargit la classification arithmétique ordinaire. Entier,
rationnel, irrationnel, algébrique et transcendant continuent de désigner des
propriétés de la valeur. HMT incorpore un deuxième axe : la manière dont une
section se prolonge, revient, conserve la mémoire et atteint une évaluation
stable. La clôture des Parties~I--III fixe les deux axes et l'application qui les relie.

\subsection{Chiffre, section, valeur et symétrie}

Soit
\[
 \mathcal S_0^{\rm surv}
 \xleftarrow{\rho_{1,0}}
 \mathcal S_1^{\rm surv}
 \xleftarrow{\rho_{2,1}}
 \cdots
\]
le système inverse d'états survivants, où chaque restriction
conserve les données typées du TPK qui appartiennent au niveau précédent, et soit
\[
 \mathcal S_\infty^{\rm surv}
 :=\varprojlim_n\mathcal S_n^{\rm surv}.
\]

\begin{definition}[Chiffre]
\label{p3-83-c44-s2:def:cifra-visible-genealogica}
Un lecteur positionnel \(L\) étant fixé, le chiffre ou bloc visible à la profondeur
\(n\) est
\[
 d_{L,n}(x_n):=\pi^{\rm vis}_{L,n}(x_n).
\]
C'est la coordonnée publiée par la présentation finie \(x_n\) ; la fibre de
\(\pi^{\rm vis}_{L,n}\) conserve les états qui partagent cette coordonnée.
\end{definition}

\begin{definition}[Nombre HMT et présentation]
\label{p3-83-c44-s2:def:numero-hmt-cierre-genealogico}
Un nombre HMT est une section compatible
\[
 \boldsymbol x=(x_n)_{n\geq0}\in\mathcal S_\infty^{\rm surv},
 \qquad
 \rho_{n+1,n}(x_{n+1})=x_n.
\]
Le couple \((\boldsymbol x,L)\) est une présentation du nombre. La section
conserve l'histoire complète ; le lecteur déclare quelles coordonnées de cette
histoire sont publiées.
\end{definition}

\begin{definition}[Valeur archimédienne]
\label{p3-83-c44-s2:def:valor-arquimediano-cierre-genealogico}
Supposons que \(L\) attribue à chaque \(x_n\) un cylindre compact
\(I_{L,n}(\boldsymbol x)\subset\mathbb R\), avec
\[
 I_{L,n+1}(\boldsymbol x)\subseteq I_{L,n}(\boldsymbol x),
 \qquad
 \operatorname{diam}I_{L,n}(\boldsymbol x)\longrightarrow0.
\]
La valeur publiée par la présentation est
\[
 \operatorname{Val}_L(\boldsymbol x)
 :=\text{l'unique élément de }
 \bigcap_{n\geq0}I_{L,n}(\boldsymbol x).
\]
\end{definition}

\begin{proposition}[Existence et unicité de l'évaluation]
\label{p3-83-c44-s2:prop:evaluacion-seccion-genealogica}
Les hypothèses précédentes déterminent une application
\[
 \operatorname{Val}_L:\mathcal S_\infty^{\rm surv}\longrightarrow\mathbb R.
\]
Chaque fibre \(\operatorname{Val}_L^{-1}(r)\) conserve les généalogies que
l'évaluation archimédienne représente par la même valeur \(r\).
\end{proposition}

\begin{proof}
La propriété d'intersection finie de la famille emboîtée de compacts
garantit une intersection non vide. Deux points distincts auraient une
distance positive majorée par tous les diamètres, en contradiction avec la
convergence de ceux-ci vers zéro.
\end{proof}

\begin{definition}[Symétrie généalogique]
\label{p3-83-c44-s2:def:simetria-genealogica-hmt}
Une symétrie du système de nombres HMT est une famille de bijections
\(g=(g_n)_{n\geq0}\) qui préserve la survie et les types, et satisfait
\[
 \rho_{n+1,n}\circ g_{n+1}
 =g_n\circ\rho_{n+1,n}
 \qquad(n\geq0).
\]
L'action
\[
 g_\infty(\boldsymbol x):=(g_n x_n)_{n\geq0}
\]
est un automorphisme de la limite inverse. Un invariant généalogique est une
fonction constante sur ses orbites.
\end{definition}

La symétrie est ainsi caractérisée comme une transformation compatible à toute
profondeur. Cette notion rend comparables la publication décimale,
l'incidence exceptionnelle et la réalisation spectrale ultérieure, en conservant le
type propre de chaque sortie.

\subsection{Conservation évolutive de l'unité}

Pour \(n\geq1\), soient
\begin{equation}
 \rho_9(n)=1+((n-1)\bmod 9),
 \qquad
 \kappa_9(n)=\frac{n-\rho_9(n)}9.
 \label{p3-83-c44-s2:eq:residuo-cociente-conservacion-unidad}
\end{equation}
Alors
\[
 n=9\kappa_9(n)+\rho_9(n),
 \qquad
 \rho_9(n)\in\{1,\ldots,9\}.
\]
La coordonnée \(\rho_9\) publie la phase nonadique et représente la classe nulle
par \(9\) ; \(\kappa_9\) conserve le nombre de tours. La racine numérique
publie la phase. La réduction modulo neuf accompagnée de son quotient publie
phase et mémoire.

La réalisation hélicoïdale associée est
\begin{equation}
 \mathcal H_9(n):=
 \left(
 \cos\frac{2\pi(\rho_9(n)-1)}9,
 \sin\frac{2\pi(\rho_9(n)-1)}9,
 \kappa_9(n)+\frac{\rho_9(n)-1}{9}
 \right).
 \label{p3-83-c44-s2:eq:helice-nonadica-conservacion-unidad}
\end{equation}

\begin{proposition}[Retour de phase et avancée de mémoire]
\label{p3-83-c44-s2:prop:retorno-fase-avance-memoria}
La réalisation \eqref{p3-83-c44-s2:eq:helice-nonadica-conservacion-unidad} satisfait
\[
 \mathcal H_9(n+9)=\mathcal H_9(n)+(0,0,1).
\]
Le retour visible \(9\to1\) clôt la phase et augmente d'une unité la mémoire
de l'état.
\end{proposition}

\begin{proof}
Les identités
\(\rho_9(n+9)=\rho_9(n)\) et
\(\kappa_9(n+9)=\kappa_9(n)+1\) fixent les coordonnées circulaires et
incrémentent exactement la hauteur.
\end{proof}

Cette identité formule la conservation évolutive : l'unité revient comme
phase et le nouvel état transporte la clôture précédente comme mémoire. La
complétion cubique réalise la même loi en dimension trois :
\begin{equation}
 10^3=(9+1)^3
 =9^3+3\cdot9^2+3\cdot9+1
 =729+243+27+1
 =729+271.
 \label{p3-83-c44-s2:eq:emrh-conservacion-evolutiva-unidad}
\end{equation}
L'intérieur possède \(729\) états ; \(243+27\) constitue la couronne
perforée ; le sommet \(+1\) complète l'unité de base mille. L'identité
\(729+270=999\) décrit l'état immédiatement antérieur à la clôture, et
\(729+271=1000\) décrit la complétion. La loi
\[
 10^d=(9+1)^d
 =\sum_{r=0}^{d}\binom dr9^{d-r}
\]
transporte la même unité entre les dimensions et conserve la généalogie de ses
strates.

\subsection{Syntaxe visible et sémantique de prolongement}

Une signature locale décrit la forme actuelle d'une branche. Sa sémantique de
survie est le système de prolongements qu'elle peut soutenir. Pour
\(x_n\in\mathcal S_n^{\rm surv}\), définissons
\[
 \operatorname{Ext}_m(x_n)
 :=\{x_m\in\mathcal S_m^{\rm surv}:\rho_{m,n}(x_m)=x_n\},
 \qquad m\geq n.
\]
Deux états qui satisfont
\[
 d_{L,n}(x_n)=d_{L,n}(y_n)
\]
peuvent posséder des arbres d'extension différents. Dans le témoin fini
\(60\to81\), deux branches ayant la même signature visible possèdent, respectivement,
zéro et douze prolongements. Le témoin établit que la continuité future
ajoute de l'information à la sortie locale. Le nombre HMT incorpore cette continuité
comme partie de l'objet.

La coinduction s'exprime par la compatibilité de tous les préfixes.
Les blocs de mille ou dix mille chiffres sont des publications tronquées d'une règle
uniformément prolongeable. La génération numérique atteint toute
profondeur demandée ; le relèvement à l'état enrichi conserve,
en outre, les hypothèses d'existence, d'unicité et de compatibilité de ses fibres.

\subsection{Classification arithmétique et classification généalogique}

Dans une base \(b\geq2\), la présentation canonique d'un réel sépare les classes
visibles suivantes :
\begin{enumerate}
\item un entier appartient à \(\mathbb Z\) et admet un mot ternaire
      équilibré fini ;
\item un rationnel possède un développement fini ou ultimement périodique ;
\item un irrationnel possède un développement non ultimement périodique ;
\item un algébrique est racine d'un polynôme non nul de \(\mathbb Q[X]\) ;
\item un transcendant appartient au complémentaire de la clôture algébrique de
      \(\mathbb Q\) dans \(\mathbb C\).
\end{enumerate}
La périodicité distingue les rationnels des irrationnels ; l'équation polynomiale
distingue les algébriques des transcendants. Ce sont des axes différents. Par exemple,
\(\varphi\) est irrationnel algébrique parce que
\(\varphi^2-\varphi-1=0\) ; la transcendance de \(e\) et \(\pi\) découle de
théorèmes classiques ; une généalogie HMT distingue, en outre, les
fonctions de propagation, de clôture et d'auto-échelle qui réalisent ces valeurs.

La même séparation régit \(e+\pi\). La somme des deux présentations
définit un objet composé et publie
\[
 \operatorname{Val}(\boldsymbol e)+
 \operatorname{Val}(\boldsymbol\pi)=e+\pi.
\]
Sa classe arithmétique constitue une question propre. La construction HMT
détermine l'application qui forme le composé et conserve les deux généalogies qui
interviennent.

\subsection{La famille opératorielle d'Euler comme grammaire de clôture}

L'équation classique d'Euler appartient à une famille quadratique unique. Soit
\(J_\sigma\) un générateur qui satisfait
\[
 J_\sigma^2=\sigma I,
 \qquad \sigma\in\{-1,0,+1\}.
\]
La séparation des puissances paires et impaires fournit
\begin{equation}
 \exp(tJ_\sigma)=C_\sigma(t)I+S_\sigma(t)J_\sigma,
 \label{p3-83-c44-s2:eq:euler-trigeometrico-cierre-numero}
\end{equation}
avec les réalisations exactes
\[
\begin{array}{ccl}
J_{-1}^2=-I&:&\exp(tJ_{-1})=\cos t\,I+\sin t\,J_{-1},\\[1mm]
J_{0}^2=0&:&\exp(tJ_{0})=I+tJ_{0},\\[1mm]
J_{+1}^2=I&:&\exp(tJ_{+1})=\cosh t\,I+\sinh t\,J_{+1}.
\end{array}
\]
Elliptique, parabolique et hyperbolique sont trois régimes d'une même
exponentielle opératorielle. La classification tritique détermine la géométrie de
l'évolution.

La branche hyperbolique contient l'auto-échelle de Fibonacci. Pour
\[
 F=\begin{pmatrix}0&1\\1&1\end{pmatrix},
 \qquad
 A=F^2=\begin{pmatrix}1&1\\1&2\end{pmatrix},
\]
on a \(\det A=1\), \(\operatorname{tr}A=3\) et
\[
 \eta=\operatorname{arcosh}\frac32=2\log\varphi.
\]
Il existe donc \(J_h^2=I\) tel que
\begin{equation}
 A=\exp(\eta J_h),
 \qquad
 A^n=\cosh(n\eta)I+\sinh(n\eta)J_h.
 \label{p3-83-c44-s2:eq:fibonacci-euler-hiperbolico-cierre-numero}
\end{equation}
La relation \(\varphi^2=\varphi+1\) exprime deux lectures du même processus :
récurrence additive locale et échelle multiplicative spectrale. Cette
compatibilité offre le modèle élémentaire de la double évaluation
somme--produit qu'APP formalise sur son support commun.

\subsection{Unification typée des géométries circulaire, complexe et hyperbolique}
\label{p3-83-c44-s2:sec:unificacion-compleja-circular-hiperbolica-numero}

La motivation géométrique initiale de HMT peut se formuler exactement comme une
architecture commune de support, de frontière et de générateurs. Le plan complexe
\(\mathbb C\) contient le cercle unité
\[
 S^1=\{e^{\mathrm i\theta}:\theta\in\mathbb R\},
\]
qui constitue la réalisation elliptique de
\eqref{p3-83-c44-s2:eq:euler-trigeometrico-cierre-numero}. Le demi-plan supérieur
\(\mathbb H=\{w\in\mathbb C:\operatorname{Im}w>0\}\) est transporté vers le disque
hyperbolique par la transformation de Cayley
\[
 \mathcal C(w)=\frac{w-\mathrm i}{w+\mathrm i},
 \qquad
 \mathcal C:\mathbb H\xrightarrow{\sim}\mathbb D,
 \qquad
 \mathbb D=\{z\in\mathbb C:|z|<1\}.
\]
La frontière idéale de \(\mathbb D\) est le même cercle \(S^1\), tandis que son
intérieur porte la métrique de Poincaré
\[
 \mathrm ds_{\mathbb D}^2
 =\frac{4\,|\mathrm dz|^2}{(1-|z|^2)^2}.
\]
La fusion consiste donc en un diagramme typé : la paramétrisation
trigonométrique publie la frontière circulaire ; la structure complexe
fournit la carte ; la métrique de Poincaré réalise l'intérieur hyperbolique.
Les trois régimes \(J^2=-I,0,+I\) incorporent en outre le seuil parabolique et
situent cette convergence géométrique au sein d'une seule famille opératorielle.

\begin{proposition}[Compatibilité euclidienne et hyperbolique du quart de tour]
\label{p3-83-c44-s2:prop:cuarto-giro-poincare-rev3}
Sur le disque unité
\[
 \mathbb D=\{z\in\mathbb C:|z|<1\},
\]
l'application
\[
 \mathcal R_{\pi/2}:\mathbb D\longrightarrow\mathbb D,
 \qquad
 \mathcal R_{\pi/2}(z)=\mathrm i z,
\]
est simultanément une rotation euclidienne et une isométrie de la métrique de
Poincaré
\[
 \mathrm ds_{\mathbb D}^{2}
 =\frac{4\,|\mathrm dz|^2}{(1-|z|^2)^2}.
\]
Son extension au bord \(S^1=\partial\mathbb D\) est le quart de tour
orienté. Dans la réalisation réelle du
\cref{thm:dos-realizaciones-fuente-cuadratica-rev3}, cette application
correspond à
\(\operatorname{Exp}((\pi/2)J_{\mathrm e})=J_{\mathrm e}\).
\end{proposition}

\begin{proof}
La multiplication par \(\mathrm i\) conserve \(|z|\) et
\(|\mathrm dz|\). Elle conserve donc aussi bien la métrique euclidienne que le
numérateur et le dénominateur de \(\mathrm ds_{\mathbb D}^{2}\). L'identité
exponentielle résulte de \(J_{\mathrm e}^2=-I\).
\end{proof}

Le plan complexe, le cercle trigonométrique et le disque hyperbolique ne sont pas
identifiés comme espaces métriques : ils partagent un support coordonné et
l'automorphisme de quart de tour après déclaration de chaque réalisation.

\subsection{Interface de frontière entre la rotation quadratique et la torsion minimale}
\label{p3-83-c44-s2:sec:interfaz-frontera-rotacion-torsion-minima-rev3}

La fonction de cette interface est de transporter une rotation intérieure déjà
construite vers une donnée de bord sans confondre l'endomorphisme quadratique
avec un défaut scalaire. Fixons un point de base
\(x_0\in\partial X\), un fibré réel
\(E_{\partial}\to\partial X\), une connexion \(\nabla_{\partial}\), un
lecteur de frontière \(q_{\partial}\) et une cotétrade discrète
\(e_{\partial}\in\Omega^1(\partial X;E_{\partial})\). Nous écrivons
\[
 \operatorname{Tors}(\partial X;E_{\partial})
 :=\Omega^2(\partial X;E_{\partial}),
 \qquad
 \operatorname{Hol}_{x_0}(\nabla_{\partial})
 \subseteq\operatorname{Aut}(E_{\partial,x_0}).
\]
Une réalisation de bord de l'interface devra construire les applications
\[
 \operatorname{Res}_{\partial}:
 \langle J_{\mathrm e}\rangle
 \longrightarrow\operatorname{Hol}_{x_0}(\nabla_{\partial}),
 \qquad
 \operatorname{Def}_{\partial}:
 \operatorname{Hol}_{x_0}(\nabla_{\partial})
 \longrightarrow\operatorname{Tors}(\partial X;E_{\partial}).
\]
La torsion déterminée par la connexion et la cotétrade a le type
\begin{equation}
 T_{\partial}
 :=d_{\nabla_{\partial}}e_{\partial}
 \in\Omega^2(\partial X;E_{\partial}).
 \label{p3-83-c44-s2:eq:torsion-borde-interfaz-euler-rev3}
\end{equation}
La réalisation physique de la torsion minimale exige en outre une fonctionnelle
\begin{equation}
 \Lambda_{\partial}:
 \operatorname{Tors}(\partial X;E_{\partial})\longrightarrow\mathbb R
 \label{p3-83-c44-s2:eq:funcional-borde-delta4-interfaz-rev3}
\end{equation}
telle que les trois applications satisfassent les conditions d'entrelacement
\begin{equation}
 \operatorname{Def}_{\partial}
 \bigl(\operatorname{Res}_{\partial}(J_{\mathrm e})\bigr)=T_{\partial},
 \qquad
 \Lambda_{\partial}(T_{\partial})=\Delta_4.
 \label{p3-83-c44-s2:eq:condiciones-borde-delta4-interfaz-rev3}
\end{equation}
Ces équations spécifient le comparateur à construire ; elles n'affirment pas
que les trois applications existent. Le scalaire \(\Delta_4\) est défini par
\cref{eq:delta4-electron-rev6}.

\begin{remark}[Condition exacte du comparateur de frontière]
La racine opératorielle \(J_{\mathrm e}^2=-I\), le quart de tour
\(\operatorname{Exp}((\pi/2)J_{\mathrm e})=J_{\mathrm e}\) et le scalaire
\(\Delta_4\) se déduisent dans leurs domaines respectifs. La chaîne de bord
précédente fixe le type des morphismes qui doivent les relier. Tant que
\(\operatorname{Res}_{\partial}\), \(\operatorname{Def}_{\partial}\) et
\(\Lambda_{\partial}\) ne sont pas construits, les conditions
\eqref{p3-83-c44-s2:eq:condiciones-borde-delta4-interfaz-rev3} définissent le comparateur
requis mais n'impliquent pas une réalisation physique ; la proximité angulaire de
ses extrémités ne remplace pas cet opérateur.
\end{remark}

APP, TRIT et TPK fournissent le support discret de cette architecture. Sur
\(X_9=(\mathbb Z/9\mathbb Z)^2\), les évaluations somme et produit induisent des
projecteurs \(P_{\Sigma,d}\) et \(P_{\Pi,d}\) qui satisfont
\[
 \|[P_{\Sigma,2},P_{\Pi,2}]\|=\frac{\sqrt2}{3},
 \qquad
 \|[P_{\Sigma,3},P_{\Pi,3}]\|=\frac{\sqrt3}{4}.
\]
Cette incompatibilité calculée est le précurseur fini de Heisenberg : deux
familles d'observables partagent un support et possèdent des décompositions
spectrales incompatibles. La forme alternée du tore,
\[
 \sigma((a,b),(c,d))=ad-bc\pmod9,
\]
s'intègre dans la représentation nonadique de Weyl--Heisenberg
\[
 X|n\rangle=|n+1\rangle,\qquad
 Z|n\rangle=\omega^n|n\rangle,\qquad
 ZX=\omega XZ,
 \qquad \omega=e^{2\pi\mathrm i/9}.
\]
Le commutateur publie la phase orientée et le groupe central résultant possède
\(9^3=729\) éléments.

Le bloc central d'APP fournit simultanément la réalisation
lorentzienne locale :
\[
 B_c=
 \begin{pmatrix}7&2\\2&7\end{pmatrix}
 =9P_++5P_-,
 \qquad
 M_c:=\frac{B_c}{\sqrt{45}}
 =\exp\!\left(\frac12\log\frac95\,R\right).
\]
Con \(\eta_{1,1}=\operatorname{diag}(1,-1)\),
\[
 M_c^{\mathsf T}\eta_{1,1}M_c=\eta_{1,1},
 \qquad
 \det M_c=1,
 \qquad
 M_c\in SO^+(1,1).
\]
Ainsi, la géométrie de Lorentz et le plan de Minkowski \(1+1\) émergent comme
réalisation spectrale exacte du bloc central. Toute élévation dimensionnelle
ultérieure doit entrelacer ce bloc local avec son nouveau codomaine. La séquence
\[
 \boxed{
 \text{APP--TRIT--TPK}
 \longrightarrow
 \begin{cases}
 \text{famille elliptique--parabolique--hyperbolique},\\
 \text{bloc de Lorentz--Minkowski }1+1,\\
 \text{représentation finie de Weyl--Heisenberg}
 \end{cases}}
\]
expose les développements géométriques et opératoriels qui précèdent la clôture
généalogique du nombre.

\subsection{Opérateurs de génération de constantes : domaines, germes et évaluations}

La fonction d'une constante est fixée au moyen de quatre données : définition
mathématique, domaine, opérateur ou lecteur et conclusion déduite. La
\cref{p3-83-c44-s2:tab:semillas-cierre-genealogico-numero} réunit les objets qui doivent
rester visibles lors de la clôture de la théorie des nombres.

\begingroup
\small
\setlength{\LTpre}{0.5\baselineskip}
\setlength{\LTpost}{0.5\baselineskip}
\begin{longtable}{@{}>{\raggedright\arraybackslash}p{.145\textwidth}
                    >{\raggedright\arraybackslash}p{.235\textwidth}
                    >{\raggedright\arraybackslash}p{.34\textwidth}
                    >{\raggedright\arraybackslash}p{.20\textwidth}@{}}
\caption{Constantes et germes dans la clôture généalogique du nombre.}
\label{p3-83-c44-s2:tab:semillas-cierre-genealogico-numero}\\
\toprule
\textbf{Objet} & \textbf{Définition ou identité} &
\textbf{Domaine et application HMT} & \textbf{Conclusion et condition}\\
\midrule
\endfirsthead
\toprule
\textbf{Objet} & \textbf{Définition ou identité} &
\textbf{Domaine et application HMT} & \textbf{Conclusion et condition}\\
\midrule
\endhead
\(\pi,e,\varphi\) &
Sections publiées par des cylindres numériques compatibles. &
La filtration
\(w_6\to w_{30}\to R_{36}\to G_9\) produit l'état terminal ; les lecteurs
séparent clôture, propagation et auto-échelle. &
Pour chaque profondeur finie prescrite, l'application produit des cylindres unitaires
compatibles ; la section infinie résulte du théorème coinductif.\\
\(\alpha_{\rm HMT}\) &
Clôture dodécaphasique avec retenue sur l'état terminal. &
La double lecture conduit à
\((P,E,\Phi,K,c_{13}=0)\) et à la sortie rationnelle de \(36\) chiffres ; la voie
\(E_{21/22}\) fournit un calcul indépendant. &
La récurrence est exacte pour l'état terminal déclaré et ne reçoit ni
\(\alpha\) ni CODATA ; la déduction minimale APP\(\to K\) exige de construire son
lecteur bidirectionnel complet.\\
\(i\) interno &
\(J^2=-I\) sur le module de Paley--Witt. &
Réalise une structure complexe finie et sépare des orientations conjuguées. &
Sur le corps de base déclaré, on a exactement \(J^2=-I\).\\
\(-1/12\) &
\(\zeta(-1)=-1/12\) et réalisations opératorielles normalisées. &
Incidence \(90\to120\), défaut dodécaphasique et fonctionnelles spectrales,
chacun avec son domaine et sa normalisation. &
Égalité exacte pour chaque application déclarée.\\
\(e+\pi\) &
Somme des présentations de \(e\) et \(\pi\). &
Objet composé de l'algèbre des présentations ; conserve les deux généalogies. &
Construction exacte ; classe arithmétique séparée.\\
Euler--Mascheroni \(\gamma\) et Stieltjes \(\gamma_n\) &
\(\gamma=\operatorname*{FP}_{s=1}\zeta(s)\) ; coefficients de Laurent. &
Décomposition résiduelle modulo trois et jets de Laurent. &
Identités exactes de la fonctionnelle déclarée.\\
Apéry \(\zeta(3)\) &
Évaluation de la zêta en \(s=3\). &
Lecteur \(\mathcal Z_3(3)\) avec intervalles rationnels de préfixe. &
Évaluation exacte ; irrationalité classique.\\
Glaisher--Kinkelin \(A_{\rm GK}\) &
\(A_{\rm GK}=\exp(1/12-\zeta'(-1))\). &
Dérivée régularisée du lecteur spectral. &
Identité exacte avec normalisation explicite.\\
Euler--Kronecker \(\gamma_{F_{-3}}\) &
Partie finie logarithmique de la zêta de Dedekind de
\(F_{-3}=\mathbb Q(\sqrt{-3})\). &
Le caractère \(\chi_{-3}\) sélectionne le corps et organise
\(L'(1,\chi_{-3})/L(1,\chi_{-3})\). &
Identité exacte du caractère déclaré.\\
Meissel--Mertens \(B_1\) &
\(B_1=\lim_{x\to\infty}(\sum_{p\le x}p^{-1}-\log\log x)\). &
Partie finie de la trace sur les modes premiers. &
Représentation opératorielle exacte avec modes fixés.\\
Catalan \(G\) et Khinchin \(K_0\) &
\(G=L(2,\chi_{-4})\) ; moyenne ergodique du système de Gauss. &
Lecteurs de caractère modulo quatre et de fraction continue. &
Théorèmes classiques incorporés avec leurs dynamiques propres.\\
Feigenbaum \((\delta,\alpha_F)\) &
Échelles de l'opérateur de renormalisation en doublement de période. &
La famille dodécaphasique produit des approximants de niveau huit. &
Le calcul détermine exactement le niveau huit ; l'universalité exige
de démontrer la convergence vers le point fixe de renormalisation.\\
Rogers--Ramanujan \(R(q)\) &
Fraction continue de niveau cinq et identité de Klein--Ramanujan pour
\(j\). &
Organise le niveau \(5\), \(A_5\), la limite
\(R(q)\to\varphi^{-1}\) et la réplicabilité de Hecke--Faber. &
Identités modulaires exactes ; la comparaison de coefficients ne couvre que la
profondeur finie calculée.\\
Horloge première \(p\) &
\(\tau_p=\operatorname{ord}_p(10)\),
\(M_p=(10^{\tau_p}-1)/p\equiv0\pmod9\) pour \(p\nmid30\). &
Les nombres premiers sont des périodes irréductibles ; les composés synchronisent des horloges
au moyen du plus petit commun multiple. &
Théorème arithmétique exact dans le domaine déclaré.\\
\bottomrule
\end{longtable}
\endgroup

Partie finie, valeur spéciale, dérivée régularisée, caractère, moyenne
ergodique, échelle de renormalisation et période première sont des applications distinctes. Leur
coexistence au sein de HMT procède d'une famille typée d'opérateurs et
d'applications d'extraction.

\subsection{Double projection d'une même généalogie}

Soit \(\mathfrak X:=X_\infty^{\rm cal}\) l'espace compact des histoires
enrichies avec jauge initiale transportée. Une jauge déclarée
\(c\in\mathscr C_{\rm exc}\) étant fixée, deux évaluations agissent sur une même section :
\[
 \mathfrak R:\mathfrak X\longrightarrow\mathbb R,
 \qquad
 \mathfrak E_c:\mathfrak X\longrightarrow\mathbf{Exc},
 \qquad
 \mathfrak E_c(x):=\mathfrak E(x,c),
\]
où \(\mathfrak R\) contracte les cylindres et \(\mathfrak E_c\) conserve les supports,
les incidences, la polarité et le transport de frontière.

\begin{theorem}[Unité causale de la double projection]
\label{p3-83-c44-s2:thm:unidad-causal-doble-proyeccion-numero}
Pour toute histoire \(\boldsymbol x\in\mathfrak X\), le couple
\[
 \mathfrak D_c(\boldsymbol x)
 :=(\mathfrak R(\boldsymbol x),\mathfrak E_c(\boldsymbol x))
\]
est une publication conjointe du même antécédent. La première coordonnée
détermine la valeur archimédienne ; la seconde détermine la classe d'incidence
exceptionnelle. Lorsque les lecteurs conjoints séparent les points,
\(\mathfrak D_c\) est injective.
\end{theorem}

\begin{proof}
Les deux composantes sont définies sur le même domaine et conservent les
applications de restriction. Si
\(\mathfrak D_c(\boldsymbol x)=\mathfrak D_c(\boldsymbol y)\), la séparation
conjointe des histoires implique \(\boldsymbol x=\boldsymbol y\).
\end{proof}

La branche archimédienne produit les cylindres et les valeurs de
\(\pi,e,\varphi\). La branche d'incidence prolonge les mêmes mots et le
même sceau vers Paley, Witt, Golay, le réseau de Leech,
\(V^\natural\), le groupe Monstre et Moonshine. La relation causale est la
préimage commune : les deux évaluations conservent des invariants différents de la
même histoire. Le réel est l'ombre contractée ; la symétrie exceptionnelle est la
mémoire d'incidence publiée par le même état.

\subsection{Clôture généalogique du concept de nombre}

Le concept de nombre est fixé par la chaîne
\[
 \boxed{
 \text{état fini}
 \longrightarrow\text{prolongement compatible}
 \longrightarrow\text{section généalogique}
 \longrightarrow
 \begin{cases}
   \text{chiffre visible},\\
   \text{valeur archimédienne},\\
   \text{incidence exceptionnelle}.
 \end{cases}}
\]
Le chiffre est une coordonnée ; la valeur est une projection ; le nombre est la
section complète ; la symétrie est l'automorphisme compatible qui conserve sa
mémoire. La conservation évolutive transporte cette identité à travers le
retour nonadique, la retenue hélicoïdale et la complétion
\(729+271=1000\). Ces distinctions étant closes, la mécanique dimensionnelle
réalise la même généalogie dans un nouveau codomaine.

\paragraph{Dépendances mathématiques.}
La limite inverse fournit la section ; le lecteur archimédien produit la
valeur ; la complétion binomiale et l'extension nonadique transportent unité et
mémoire ; les lecteurs complémentaires classent les sorties ; et la double
projection applique au même antécédent les applications archimédienne et exceptionnelle.
La chaîne Witt--Golay--Leech--Moonshine exige, à chaque étape,
l'entrelaceur d'incidence correspondant.
